
\Data\ID{1003024601}
\Updated{2020-07-19 02:07:04}
\Level{A}
\SubArea{Combinatorics}
\LANGen{
\Question{
Assume the password for the safe deposit box consists of four different letters from the set \( \{A;B;C;D;E;F;G;H\} \) and four different numbers from the set \( \{1;2;3;4;5;6;7\} \). How many different passwords are there?}
{\( \binom84 \cdot \binom74 \cdot 8! = 98\,784\,000 \)}{\( \frac{8!}{4!}\cdot\frac{7!}{3!}\cdot8!=56\,899\,584\,000 \)}{\( \left(\frac{8!}{4!}+\frac{7!}{3!}\right)\cdot8! = 101\,606\,400  \)}{\( \left(\binom84+\binom74\right)\cdot8!=4\,233\,600 \)}{}{}}
\LANGpl{
\Question{
Załóżmy, że hasło do sejfu składa się z czterech różnych liter ze zbioru \( \{A;B;C;D;E;F;G;H\} \) i czterech różnych cyfr ze zbioru \ \( \{1;2;3;4;5;6;7\} \). Ile jest różnych haseł?}
{\( \binom84 \cdot \binom74 \cdot 8! = 98\,784\,000 \)}{\( \frac{8!}{4!}\cdot\frac{7!}{3!}\cdot8!=56\,899\,584\,000 \)}{\( \left(\frac{8!}{4!}+\frac{7!}{3!}\right)\cdot8! = 101\,606\,400  \)}{\( \left(\binom84+\binom74\right)\cdot8!=4\,233\,600 \)}{}{}}
\LANGcs{
\Question{
Heslem v trezoru je skupina libovolně uspořádaných písmen a číslic skládajících se ze čtyř různých písmen z množiny \( \{A;B;C;D;E;F;G;H\} \) a čtyř různých číslic z množiny \( \{1;2;3;4;5;6;7\} \). Kolik různých hesel se dá v trezoru nastavit?}
{\( \binom84 \cdot \binom74 \cdot 8! = 98\,784\,000 \)}{\( \frac{8!}{4!}\cdot\frac{7!}{3!}\cdot8!=56\,899\,584\,000 \)}{\( \left(\frac{8!}{4!}+\frac{7!}{3!}\right)\cdot8! = 101\,606\,400  \)}{\( \left(\binom84+\binom74\right)\cdot8!=4\,233\,600 \)}{}{}}
\LANGsk{
\Question{
Heslom v trezore je skupina ľubovoľne usporiadaných písmen a číslic pozostávajúca zo štyroch rôznych písmen z množiny \( \{A;B;C;D;E;F;G;H\} \) a štyroch rôznych číslic z množiny \( \{1;2;3;4;5;6;7\} \). Koľko rôznych hesiel sa dá v trezore nastaviť?}
{\( \binom84 \cdot \binom74 \cdot 8! = 98\,784\,000 \)}{\( \frac{8!}{4!}\cdot\frac{7!}{3!}\cdot8!=56\,899\,584\,000 \)}{\( \left(\frac{8!}{4!}+\frac{7!}{3!}\right)\cdot8! = 101\,606\,400  \)}{\( \left(\binom84+\binom74\right)\cdot8!=4\,233\,600 \)}{}{}}
\LANGes{
\Question{
Supón que la contraseña para una caja de seguridad consta de cuatro letras diferentes del conjunto \( \{A;B;C;D;E;F;G;H\} \) y cuatro números diferentes del conjunto \( \{1;2;3;4;5;6;7\} \). ¿Cuántas contraseñas diferentes hay?}
{\( \binom84 \cdot \binom74 \cdot 8! = 98\,784\,000 \)}{\( \frac{8!}{4!}\cdot\frac{7!}{3!}\cdot8!=56\,899\,584\,000 \)}{\( \left(\frac{8!}{4!}+\frac{7!}{3!}\right)\cdot8! = 101\,606\,400  \)}{\( \left(\binom84+\binom74\right)\cdot8!=4\,233\,600 \)}{}{}}
\End

\Data\ID{1003024602}
\Updated{2020-07-19 03:07:57}
\Level{A}
\SubArea{Combinatorics}
\LANGen{
\Question{
Suppose six pupils are sitting on six chairs in a row. Among them there are twins. In how many ways can the pupils be seated in the row so that the twins do not sit next to each other?}
{\( 6! -2\cdot5!=480 \)}{\( \frac{6!}2=360 \)}{\( 2\cdot5!=240 \)}{\( 6! -5!=600 \)}{}{}}
\LANGpl{
\Question{
Załóżmy, że sześciu uczniów siedzi na sześciu krzesłach ustawionych w  rzędzie. Wśród nich są bliźnięta. Na ile sposobów uczniowie mogą usiąść w rzędzie, aby bliźnięta nie siedziały obok siebie?}
{\( 6! -2\cdot5!=480 \)}{\( \frac{6!}2=360 \)}{\( 2\cdot5!=240 \)}{\( 6! -5!=600 \)}{}{}}
\LANGcs{
\Question{
V lavici na 6 židlích má sedět 6 žáků, mezi kterými jsou dvojčata. Kolika způsoby můžeme žáky v lavici posadit tak, aby dvojčata neseděla vedle sebe?}
{\( 6! -2\cdot5!=480 \)}{\( \frac{6!}2=360 \)}{\( 2\cdot5!=240 \)}{\( 6! -5!=600 \)}{}{}}
\LANGsk{
\Question{
V lavici na 6 stoličkách má sedieť 6 žiakov, medzi ktorými sú dvojčatá. Koľkými spôsobmi môžeme žiakov v lavici posadiť tak, aby dvojčatá nesedeli vedľa seba?}
{\( 6! -2\cdot5!=480 \)}{\( \frac{6!}2=360 \)}{\( 2\cdot5!=240 \)}{\( 6! -5!=600 \)}{}{}}
\LANGes{
\Question{
Supón que seis alumnos están sentados en seis sillas en una fila. Entre ellos hay gemelos. ¿De cuántas maneras pueden sentarse los alumnos en la fila para que los gemelos no se sienten uno al lado del otro?}
{\( 6! -2\cdot5!=480 \)}{\( \frac{6!}2=360 \)}{\( 2\cdot5!=240 \)}{\( 6! -5!=600 \)}{}{}}
\End

\Data\ID{1003024606}
\Updated{2021-04-25 06:04:24}
\Level{A}
\SubArea{Combinatorics}
\LANGen{
\Question{
Each payment card has its numeric four-digit PIN code. How many different PIN codes can be selected, if only a code with different numbers may be used?}
{\( \frac{10!}{6!} = 5\:040 \)}{\( \frac{10!}{4!} = 151\:200 \)}{\( \frac{10!}{6!\cdot4!} = 210 \)}{\( 10^4 = 10\:000 \)}{}{}}
\LANGpl{
\Question{
Każda karta płatnicza ma swój czterocyfrowy numer PIN. Ile różnych kodów PIN można utworzyć, jeśli kod musi mieć  różne cyfry?}
{\( \frac{10!}{6!} = 5\:040 \)}{\( \frac{10!}{4!} = 151\:200 \)}{\( \frac{10!}{6!\cdot4!} = 210 \)}{\( 10^4 = 10\:000 \)}{}{}}
\LANGcs{
\Question{
Každá platební karta má svůj číselný čtyřciferný PIN kód. Kolik různých PIN kódů je možné zvolit, když se kvůli bezpečnosti použije jen kód s různými číslicemi?}
{\( \frac{10!}{6!} = 5\:040 \)}{\( \frac{10!}{4!} = 151\:200 \)}{\( \frac{10!}{6!\cdot4!} = 210 \)}{\( 10^4 = 10\:000 \)}{}{}}
\LANGsk{
\Question{
Každá platobná karta má svoj číselný štvorciferný PIN kód. Koľko rôznych PIN kódov možno zvoliť, ak sa kvôli bezpečnosti použije len kód s rôznymi číslicami?}
{\( \frac{10!}{6!} = 5\:040 \)}{\( \frac{10!}{4!} = 151\:200 \)}{\( \frac{10!}{6!\cdot4!} = 210 \)}{\( 10^4 = 10\:000 \)}{}{}}
\LANGes{
\Question{
Cada tarjeta de crédito tiene un PIN de cuatro dígitos. ¿Cuántos códigos diferentes de cuatro dígitos  se pueden seleccionar si solo se puede usar un código con números diferentes?}
{\( \frac{10!}{6!} = 5\:040 \)}{\( \frac{10!}{4!} = 151\:200 \)}{\( \frac{10!}{6!\cdot4!} = 210 \)}{\( 10^4 = 10\:000 \)}{}{}}
\End

\Data\ID{1003024607}
\Updated{2020-07-19 03:07:15}
\Level{A}
\SubArea{Combinatorics}
\LANGen{
\Question{
On the shelf, there should be three blue cups, three red cups, two yellow cups and two green cups arranged in a row from the left to the right. The cups of the same color are not mutually distinguishable. How many arrangements of these cups are possible?}
{\( \frac{10!}{(2!)^2\cdot(3!)^2}=25\:200 \)}{\( \frac{10!}{4\cdot6!}=1\:260 \)}{\( \frac{10!}{2\cdot2!\cdot3!}=151\:200 \)}{\( \frac{10!}{4\cdot2!\cdot3!}=75\:600 \)}{}{}}
\LANGpl{
\Question{
Na półce znajdują się trzy niebieskie kubki, trzy czerwone kubki, dwa żółte i dwa zielone kubki  rozmieszczone w rzędzie od lewej do prawej. Kubki tego samego koloru nie różnią się od siebie. Ile jest możliwych ułożeń tych kubków?}
{\( \frac{10!}{(2!)^2\cdot(3!)^2}=25\:200 \)}{\( \frac{10!}{4\cdot6!}=1\:260 \)}{\( \frac{10!}{2\cdot2!\cdot3!}=151\:200 \)}{\( \frac{10!}{4\cdot2!\cdot3!}=75\:600 \)}{}{}}
\LANGcs{
\Question{
Na poličce je třeba v řadě (zleva doprava) uložit tři modré, tři červené, dva žluté a dva zelené hrnečky, přičemž hrnečky stejné barvy jsou navzájem nerozlišitelné. Kolik je možných různých uložení těchto hrnečků?}
{\( \frac{10!}{(2!)^2\cdot(3!)^2}=25\:200 \)}{\( \frac{10!}{4\cdot6!}=1\:260 \)}{\( \frac{10!}{2\cdot2!\cdot3!}=151\:200 \)}{\( \frac{10!}{4\cdot2!\cdot3!}=75\:600 \)}{}{}}
\LANGsk{
\Question{
Na poličke treba v rade (zľava doprava) uložiť tri modré, tri červené, dva žlté a dva zelené hrnčeky, pričom hrnčeky rovnakej farby sú navzájom nerozlíšiteľné. Koľko je možných rôznych uložení týchto hrnčekov?}
{\( \frac{10!}{(2!)^2\cdot(3!)^2}=25\:200 \)}{\( \frac{10!}{4\cdot6!}=1\:260 \)}{\( \frac{10!}{2\cdot2!\cdot3!}=151\:200 \)}{\( \frac{10!}{4\cdot2!\cdot3!}=75\:600 \)}{}{}}
\LANGes{
\Question{
En la estantería deberían encontrarse tres tazas azules, tres tazas rojas, dos tazas amarillas y dos tazas verdes arregladas en una fila de izquierda a derecha. Las tazas del mismo color no se pueden distinguir entre sí. ¿Cuántos arreglos de estas tazas son posibles?}
{\( \frac{10!}{(2!)^2\cdot(3!)^2}=25\:200 \)}{\( \frac{10!}{4\cdot6!}=1\:260 \)}{\( \frac{10!}{2\cdot2!\cdot3!}=151\:200 \)}{\( \frac{10!}{4\cdot2!\cdot3!}=75\:600 \)}{}{}}
\End

\Data\ID{1003024608}
\Updated{2020-07-19 03:07:04}
\Level{A}
\SubArea{Combinatorics}
\LANGen{
\Question{
How many ways are there to arrange the letters in the czech word KOMBINATORIKA?}
{\( \frac{13!}{(2!)^4}=389\:188\:800 \)}{\( \frac{13!}{4\cdot2!}=778\:377\:600 \)}{\( \frac{13!}{2!}=3\:113\:510\:400 \)}{\( 13!-4\cdot2!=6\:227\:020\:792 \)}{}{}}
\LANGpl{
\Question{
Na ile sposobów można ułożyć litery w czeskim słowie KOMBINATORIKA?}
{\( \frac{13!}{(2!)^4}=389\:188\:800 \)}{\( \frac{13!}{4\cdot2!}=778\:377\:600 \)}{\( \frac{13!}{2!}=3\:113\:510\:400 \)}{\( 13!-4\cdot2!=6\:227\:020\:792 \)}{}{}}
\LANGcs{
\Question{
Kolik různých anagramů (přesmyček, které ale nemusí mít žádný význam) je možné sestavit ze všech písmen slova KOMBINATORIKA?}
{\( \frac{13!}{(2!)^4}=389\:188\:800 \)}{\( \frac{13!}{4\cdot2!}=778\:377\:600 \)}{\( \frac{13!}{2!}=3\:113\:510\:400 \)}{\( 13!-4\cdot2!=6\:227\:020\:792 \)}{}{}}
\LANGsk{
\Question{
Koľko rôznych anagramov (prešmyčiek, ktoré ale nemusia mať žiadny význam) možno zostaviť zo všetkých písmen slova KOMBINATORIKA?}
{\( \frac{13!}{(2!)^4}=389\:188\:800 \)}{\( \frac{13!}{4\cdot2!}=778\:377\:600 \)}{\( \frac{13!}{2!}=3\:113\:510\:400 \)}{\( 13!-4\cdot2!=6\:227\:020\:792 \)}{}{}}
\LANGes{
\Question{
¿De cuántas maneras se pueden ordenar las letras en la palabra checa "KOMBINATORIKA"?}
{\( \frac{13!}{(2!)^4}=389\:188\:800 \)}{\( \frac{13!}{4\cdot2!}=778\:377\:600 \)}{\( \frac{13!}{2!}=3\:113\:510\:400 \)}{\( 13!-4\cdot2!=6\:227\:020\:792 \)}{}{}}
\End

\Data\ID{1003024609}
\Updated{2021-04-25 06:04:43}
\Level{A}
\SubArea{Combinatorics}
\LANGen{
\Question{
How many \( 8 \)-digit numbers are there in which the digit sum is \( 3 \)?}
{\( 36 \)}{\( 113 \)}{\( 112 \)}{\( 35 \)}{}{}}
\LANGpl{
\Question{
Ile jest \( 8 \)-cyfrowych liczb, w których suma cyfr to \( 3 \)?}
{\( 36 \)}{\( 113 \)}{\( 112 \)}{\( 35 \)}{}{}}
\LANGcs{
\Question{
Kolik je všech \( 8 \)-ciferných čísel, ve kterých je ciferný součet \( 3 \)?}
{\( 36 \)}{\( 113 \)}{\( 112 \)}{\( 35 \)}{}{}}
\LANGsk{
\Question{
Koľko je všetkých \( 8 \)-ciferných čísel, v ktorých ciferný súčet je \( 3 \)?}
{\( 36 \)}{\( 113 \)}{\( 112 \)}{\( 35 \)}{}{}}
\LANGes{
\Question{
¿Cuántos números de  \( 8 \) dígitos existen cuya suma de sus propios dígitos es  \( 3 \)?}
{\( 36 \)}{\( 113 \)}{\( 112 \)}{\( 35 \)}{}{}}
\End

\Data\ID{1003024610}
\Updated{2021-04-25 06:04:21}
\Level{A}
\SubArea{Combinatorics}
\LANGen{
\Question{
In a high-speed train set, there should be the following cars included: \( 3 \) first class cars, \( 5 \) second class cars, \( 2 \) sleeping cars, \( 1 \) dining car, and \( 2 \) luggage cars. How many ways are there to arrange the cars in this high-speed train set?}
{\( \frac{13!}{(2!)^2\cdot3!\cdot5!}=2\:162\:160 \)}{\( \frac{13!}{(2!)^2+3!+5!}=47\:900\:160 \)}{\( 13!-(2!)^2\cdot3!\cdot5!=6\:227\:017\:920 \)}{\( 13!-\left|(2!)^2+3!+5!\right|=6\:227\:020\:670 \)}{}{}}
\LANGpl{
\Question{
W pociągu ekspresowym powinny być następujące wagony: \( 3 \) pierwszej klasy, \( 5 \) drugiej klasy, \( 2 \) sypialniane, \( 1 \) wagon gastronomiczny oraz  \( 2 \) wagony bagażowe. Ile jest sposobów ułożenia wagonów w pociągu?}
{\( \frac{13!}{(2!)^2\cdot3!\cdot5!}=2\:162\:160 \)}{\( \frac{13!}{(2!)^2+3!+5!}=47\:900\:160 \)}{\( 13!-(2!)^2\cdot3!\cdot5!=6\:227\:017\:920 \)}{\( 13!-\left|(2!)^2+3!+5!\right|=6\:227\:020\:670 \)}{}{}}
\LANGcs{
\Question{
Do rychlíkové soupravy mají být zařazené tyto druhy vagónů: \( 3 \) vozy 1. třídy, \( 5 \) vozů 2. třídy, \( 2 \) lůžkové,  \( 1 \) jídelní a  \( 2 \) zavazadlové vozy. Kolika různými způsoby z nich můžeme sestavit rychlíkovou soupravu?}
{\( \frac{13!}{(2!)^2\cdot3!\cdot5!}=2\:162\:160 \)}{\( \frac{13!}{(2!)^2+3!+5!}=47\:900\:160 \)}{\( 13!-(2!)^2\cdot3!\cdot5!=6\:227\:017\:920 \)}{\( 13!-\left|(2!)^2+3!+5!\right|=6\:227\:020\:670 \)}{}{}}
\LANGsk{
\Question{
Do rýchlikovej súpravy majú byť zaradené tieto druhy vagónov: \( 3 \) vozy 1. triedy, \( 5 \) vozov 2. triedy, \( 2 \) lôžkové,  \( 1 \) jedálenský a  \( 2 \) batožinové vozy. Koľkými rôznymi spôsobmi z nich môžeme zostaviť rýchlikovú súpravu?}
{\( \frac{13!}{(2!)^2\cdot3!\cdot5!}=2\:162\:160 \)}{\( \frac{13!}{(2!)^2+3!+5!}=47\:900\:160 \)}{\( 13!-(2!)^2\cdot3!\cdot5!=6\:227\:017\:920 \)}{\( 13!-\left|(2!)^2+3!+5!\right|=6\:227\:020\:670 \)}{}{}}
\LANGes{
\Question{
En los trenes de alta velocidad debe haber los siguientes tipos de vagones incluidos: \(3 \) vagones de primera clase, \(5 \) vagones de segunda clase, \(2 \) vagones para dormir, \(1 \) vagón comedor, y \(2 \) vagones de equipaje. ¿De cuántas maneras se pueden organizar los vagones en los trenes de alta velocidad?}
{\( \frac{13!}{(2!)^2\cdot3!\cdot5!}=2\:162\:160 \)}{\( \frac{13!}{(2!)^2+3!+5!}=47\:900\:160 \)}{\( 13!-(2!)^2\cdot3!\cdot5!=6\:227\:017\:920 \)}{\( 13!-\left|(2!)^2+3!+5!\right|=6\:227\:020\:670 \)}{}{}}
\End

\Data\ID{1003024611}
\Updated{2021-04-25 06:04:43}
\Level{A}
\SubArea{Combinatorics}
\LANGen{
\Question{
On a lock of a safe deposit box a ten-digit code can be set. The code can consist only of four \( 1 \)s, three \( 2 \)s, two \( 3 \)s, and one \( 4 \). How many ways are there to set the code?}
{\( \frac{10!}{4!\cdot3!\cdot2!} = 12\:600 \)}{\( \frac{10!}{4!+3!+2!}=113\:400 \)}{\( 10!-4!\cdot3!\cdot5!=3\:628\:512 \)}{\( 10! = 3\:628\:800 \)}{}{}}
\LANGpl{
\Question{
W sejfie można ustawić dziesięciocyfrowy kod. Kod może składać się tylko z czterech \( 1 \), trzech \( 2 \), dwóch \( 3 \)  oraz jednej  \( 4 \). Ile rodzajów kodu można utworzyć?}
{\( \frac{10!}{4!\cdot3!\cdot2!} = 12\:600 \)}{\( \frac{10!}{4!+3!+2!}=113\:400 \)}{\( 10!-4!\cdot3!\cdot5!=3\:628\:512 \)}{\( 10! = 3\:628\:800 \)}{}{}}
\LANGcs{
\Question{
Na zámku trezoru je třeba nastavit deseticiferný kód, který se skládá jen ze čtyř jednotek, třech dvojek, dvou trojek a jedné čtyřky. Kolika různými způsoby to je možné udělat?}
{\( \frac{10!}{4!\cdot3!\cdot2!} = 12\:600 \)}{\( \frac{10!}{4!+3!+2!}=113\:400 \)}{\( 10!-4!\cdot3!\cdot5!=3\:628\:512 \)}{\( 10! = 3\:628\:800 \)}{}{}}
\LANGsk{
\Question{
Na zámku trezora treba nastaviť desaťciferný kód, ktorý pozostáva len zo štyroch jednotiek, troch dvojok, dvoch trojok a jednej štvorky. Koľkými rôznymi spôsobmi to možno urobiť?}
{\( \frac{10!}{4!\cdot3!\cdot2!} = 12\:600 \)}{\( \frac{10!}{4!+3!+2!}=113\:400 \)}{\( 10!-4!\cdot3!\cdot5!=3\:628\:512 \)}{\( 10! = 3\:628\:800 \)}{}{}}
\LANGes{
\Question{
En la cerradura de una caja de seguridad se puede configurar un código de diez dígitos. El código puede consistir solo en cuatro \(1 \) s, tres \(2 \) s, dos \(3 \) s y un \(4 \). ¿De cuántas maneras se puede configurar el código?}
{\( \frac{10!}{4!\cdot3!\cdot2!} = 12\:600 \)}{\( \frac{10!}{4!+3!+2!}=113\:400 \)}{\( 10!-4!\cdot3!\cdot5!=3\:628\:512 \)}{\( 10! = 3\:628\:800 \)}{}{}}
\End

\Data\ID{9000139301}
\Updated{2021-04-25 06:04:30}
\Level{A}
\SubArea{Combinatorics}
\LANGen{
\Question{
Five different pieces of fruit are there on the plate. In how many ways it is possible
to distribute them into five children, if each child should have one piece?}
{\(120\)}{\(100\)}{\(140\)}{\(160\)}{}{}}
\LANGpl{
\Question{
Na talerzu jest pięć kawałków owoców (jabłko, gruszka, kiwi, banan i pomarańcza). Ile jest sposobów podzielenia owoców między pięcioro dzieci  tak, aby każde dziecko otrzymało jeden owoc?}
{\(120\)}{\(100\)}{\(140\)}{\(160\)}{}{}}
\LANGcs{
\Question{
Na talíři je pět kusů ovoce (jablko, hruška, kiwi, banán a pomeranč).
Kolika způsoby je možné ovoce rozdělit mezi pět dětí tak, aby každé
dítě dostalo jeden kus ovoce?}
{\(120\)}{\(100\)}{\(140\)}{\(160\)}{}{}}
\LANGsk{
\Question{
Na tanieri je 5 rôznych kusov ovocia. Koľkými spôsobmi je možné ovocie rozdeliť medzi 5 detí tak, aby každé dieťa dostalo jeden kus ovocia?}
{\(120\)}{\(100\)}{\(140\)}{\(160\)}{}{}}
\LANGes{
\Question{
En un plato hay cinco frutas diferentes. ¿De cuántas formas es posible
distribuirlas entre cinco niños, si cada niño debe tener una fruta?}
{\(120\)}{\(100\)}{\(140\)}{\(160\)}{}{}}
\End

\Data\ID{9000139302}
\Updated{2021-04-25 06:04:13}
\Level{A}
\SubArea{Combinatorics}
\LANGen{
\Question{
The phone number contains nine digits. A witness does not remember
the full number, but he remembers that the phone number starts by
\(728\), ends
by \(01\) 
and there is no repeating digit in the number. How many phone numbers meet these
conditions?}
{\(120\)}{\(320\)}{\(520\)}{\(720\)}{}{}}
\LANGpl{
\Question{
Numer telefonu zawiera dziewięć cyfr. Świadek nie pamięta
pełnego numeru, ale pamięta, że numer telefonu zaczyna się od
\(728\), a kończy 
na  \(01\) 
ponadto żadna z cyfr nie powtarza się. Ile numerów telefonów spełnia te warunki?}
{\(120\)}{\(320\)}{\(520\)}{\(720\)}{}{}}
\LANGcs{
\Question{
O devítimístném telefonním čísle jsme si zapamatovali jen to, že začíná
trojčíslím \(728\),
každá číslice se vyskytuje jen jednou a číslo má na konci dvojčíslí
\(01\). Kolik
telefonních čísel odpovídá popisu?}
{\(120\)}{\(320\)}{\(520\)}{\(720\)}{}{}}
\LANGsk{
\Question{
Telefónne číslo obsahuje deväť číslic. Zapamätali sme si len to, že začína
trojčíslom 
\(728\), končí dvojčíslom
 \(01\) 
a každá číslica sa vyskytuje len raz. Koľko telefónnych čísel zodpovedá popisu?}
{\(120\)}{\(320\)}{\(520\)}{\(720\)}{}{}}
\LANGes{
\Question{
Un número de teléfono contiene nueve dígitos. Una persona no recuerda
el número completo, pero recuerda que el número de teléfono empieza por
\(728\), termina
por \(01\) 
y que no hay ningún dígito repetido en el número. ¿Cuántos números de teléfono cumplen estas
condiciones?}
{\(120\)}{\(320\)}{\(520\)}{\(720\)}{}{}}
\End

\Data\ID{9000139303}
\Updated{2020-08-18 06:08:55}
\Level{A}
\SubArea{Combinatorics}
\LANGen{
\Question{
A DJ's playlist contains \(18\) 
songs. In this list there are \(7\) 
rap songs, \(5\) 
oldies and \(6\) 
rock songs. The opening part should consist of one rap song, two oldies and one rock
song. The order of the songs does not matter. Find the number of possible ways how
to put the opening together.}
{\(420\)}{\(120\)}{\(320\)}{\(520\)}{}{}}
\LANGpl{
\Question{
Lista odtwarzania DJ zawiera \(18\) 
piosenek, z czego  \(7\) to piosenki
rap, \(5\) 
house  i \(6\) 
rock. Część na otwarcie powinna zawierać piosenkę rap,  dwie typu house, i jedną typu  rock.
 Kolejność utworów nie ma znaczenia.  Na ile możliwych sposobów można ułożyć listę na otwarcie?}
{\(420\)}{\(120\)}{\(320\)}{\(520\)}{}{}}
\LANGcs{
\Question{
Klubový DJ má na hodinové představení na playlistu nachystáno
\(18\) různých
písniček, z toho \(7\) 
je z kategorie techno, \(5\) 
oldies a \(6\) 
house. Na první úvodní část chce vybrat jednu písničku techno, dvě
oldies a jednu house. Kolik je možností sestavení úvodního playlistu,
nezáleží-li nám na pořadí vybraných písniček?}
{\(420\)}{\(120\)}{\(320\)}{\(520\)}{}{}}
\LANGsk{
\Question{
Klubový DJ má na svojom zozname skladieb nachystaných \(18\) rôznych
piesní, z toho \(7\) je z kategórie
rap, \(5\) 
oldies a \(6\) 
rock. Úvodná časť predstavenia by mala obsahovať jednú pieseň rap, dve oldies a jednu rockovú pieseň. Koľko je možností zostavenia úvodného zoznamu skladieb, ak  nám nezáleží na poradí vybraných pesničiek?}
{\(420\)}{\(120\)}{\(320\)}{\(520\)}{}{}}
\LANGes{
\Question{
La lista de reproducción de un DJ contiene  \(18\) 
canciones. En esta lista hay \(7\) 
canciones de rap, \(5\) 
canciones clásicas y \(6\) 
canciones de rock. El incicio de la lista de reproducción debe constar de una canción de rap, dos canciones clásicas y una de rock. El orden de las canciones no importa. Calcula el número de formas posibles de cómo iniciar la lista de reproducción.}
{\(420\)}{\(120\)}{\(320\)}{\(520\)}{}{}}
\End

\Data\ID{9000139304}
\Updated{2020-08-18 06:08:42}
\Level{A}
\SubArea{Combinatorics}
\LANGen{
\Question{
Find the number of possibilities how to choose a pair of timekeepers for a sport event if
there are \(50\) 
candidates available for this job.}
{\(\frac{50!}
{2!\; 48!} = 1\:225\)}{\(\frac{50}
 {2} = 25\)}{\(50^{2} = 2\:500\)}{\(\frac{50!}
{48!} = 2\:450\)}{}{}}
\LANGpl{
\Question{
Znajdź liczbę możliwości, wybrania  pary piłkarzy do imprezy sportowej, jeśli
jest \(50\) 
kandydatów do wyboru.}
{\(\frac{50!}
{2!\; 48!} = 1\:225\)}{\(\frac{50}
 {2}= 25 \)}{\(50^{2}= 2\:500\)}{\(\frac{50!}
{48!}=2\:450\)}{}{}}
\LANGcs{
\Question{
Kolika způsoby je možné vybrat ze skupiny
\(50\) 
pořadatelů dvojici, která bude měřit závodníkům čas?}
{\(\frac{50!}
{2!\; 48!} = 1\:225\)}{\(\frac{50}
 {2} = 25\)}{\(50^{2} = 2\:500\)}{\(\frac{50!}
{48!} = 2\:450\)}{}{}}
\LANGsk{
\Question{
Koľkými spôsobmi je možné vybrať zo skupiny \(50\) usporiadateľov dvoch, ktorí budú merať pretekárom čas?}
{\(\frac{50!}
{2!\; 48!}=1\:250\)}{\(\frac{50}
 {2} =25\)}{\(50^{2}=2\:500\)}{\(\frac{50!}
{48!}=2\:450\)}{}{}}
\LANGes{
\Question{
Halla el número de posibilidades de cómo elegir una pareja de cronometradores para un evento deportivo si
hay \(50\)
candidatos disponibles para este trabajo.}
{\(\frac{50!}
{2!\; 48!} = 1\:225\)}{\(\frac{50}
 {2} = 25\)}{\(50^{2} = 2\:500\)}{\(\frac{50!}
{48!} = 2\:450\)}{}{}}
\End

\Data\ID{9000139305}
\Updated{2020-08-18 07:08:27}
\Level{A}
\SubArea{Combinatorics}
\LANGen{
\Question{
There are five rooms with three beds and one room with five beds in a hotel. A group of
\(20\) 
people booked rooms at this hotel. Determine the number of possible choices for the
people to the five-bed room.}
{\(\frac{20!}
{5!\; 15!}=15\:504\)}{\(20\cdot 3\cdot 5=300\)}{\(\frac{20!}
{3!\; 5!}=3\:379\:030\:566\:912\:000\)}{\(20^{5}=3\:200\:000\)}{}{}}
\LANGpl{
\Question{
Jest pięć pokoi z trzema łóżkami i jeden pokój z pięcioma łóżkami w hotelu. Grupa
\(20\) osób dokonała rezerwacji w  hotelu. Określ liczbę możliwości zakwaterowania tych osób w pokoju pięcioosobowym.}
{\(\frac{20!}
{5!\; 15!}=15\:504\)}{\(20\cdot 3\cdot 5=300\)}{\(\frac{20!}
{3!\; 5!}=3\:379\:030\:566\:912\:000\)}{\(20^{5}=3\:200\:000\)}{}{}}
\LANGcs{
\Question{
Skupina \(20\) 
studentů se má ubytovat v penzionu. Recepční má k dispozici
\(5\) třílůžkových
pokojů a \(1\) 
pětilůžkový. Kolika způsoby je možné vybrat pět studentů, kteří
budou ubytování v pětilůžkovém pokoji?}
{\(\frac{20!}
{5!\; 15!}=15\:504\)}{\(20\cdot 3\cdot 5=300\)}{\(\frac{20!}
{3!\; 5!}=3\:379\:030\:566\:912\:000\)}{\(20^{5}=3\:200\:000\)}{}{}}
\LANGsk{
\Question{
V hoteli je päť izieb s tromi lôžkami a jedna izba s piatimi lôžkami. Skupina
\(20\) 
ľudí sa chce ubytovať v tomto hoteli. Koľkými spôsobmi je možné vybrať päť ľudí, ktorí budú ubytovaní v izbe s piatimi lôžkami?}
{\(\frac{20!}
{5!\; 15!}=15\:504\)}{\(20\cdot 3\cdot 5=300\)}{\(\frac{20!}
{3!\; 5!}=3\:379\:030\:566\:912\:000\)}{\(20^{5}=3\:200\:000\)}{}{}}
\LANGes{
\Question{
Hay cinco habitaciones con tres camas y una habitación con cinco camas en un hotel. Un grupo de
\(20\) 
personas hicieron una reserva de habitaciones en este hotel. ¿De cuántas formas pueden ser elegidos cinco estudiantes que 
serán alojados en la habitación de cinco camas?}
{\(\frac{20!}
{5!\; 15!}=15\:504\)}{\(20\cdot 3\cdot 5=300\)}{\(\frac{20!}
{3!\; 5!}=3\:379\:030\:566\:912\:000\)}{\(20^{5}=3\:200\:000\)}{}{}}
\End

\Data\ID{9000139306}
\Updated{2021-04-25 06:04:31}
\Level{A}
\SubArea{Combinatorics}
\LANGen{
\Question{
How many different ice cream bowls can be prepared from
\(8\) kinds
of an ice cream if the bowl contains just three different kinds of ice cream?}
{\(\frac{8!}
{3!\; 5!}=56\)}{\(\frac{8!}
{2!\; 5!}=168\)}{\(\frac{8!}
{5!}=336\)}{\(8!=40\:320\)}{}{}}
\LANGpl{
\Question{
W cukierni jest \(8\)  rodzajów lodów. Deser lodowy może zawierać trzy różne rodzaje lodów. Wskaż na ile możliwych sposobów może zostać przygotowany deser lodowy.}
{\(\frac{8!}
{3!\; 5!}=56\)}{\(\frac{8!}
{2!\; 5!}=168\)}{\(\frac{8!}
{5!}=336\)}{\(8!=40\:320\)}{}{}}
\LANGcs{
\Question{
V cukrárně mají \(8\) 
druhů zmrzliny. Do poháru si chci objednat tři různé druhy. Kolik existuje
možností výběru?}
{\(\frac{8!}
{3!\; 5!}=56\)}{\(\frac{8!}
{2!\; 5!}=168\)}{\(\frac{8!}
{5!}=336\)}{\(8!=40\:320\)}{}{}}
\LANGsk{
\Question{
Koľko rozdielnych zmrzlinových pohárov sa dá pripraviť z 
\(8\) druhov
zmrzliny, ak pohár obsahuje len tri rôzne druhy zmrzliny?}
{\(\frac{8!}
{3!\; 5!}=56\)}{\(\frac{8!}
{2!\; 5!}=168\)}{\(\frac{8!}
{5!}=336\)}{\(8!=40\:320\)}{}{}}
\LANGes{
\Question{
¿Cuántas bolas diferentes de helado se pueden preparar con
\(8\) tipos
de un helado si la bola contiene solo tres tipos diferentes del helado?}
{\(\frac{8!}
{3!\; 5!}=56\)}{\(\frac{8!}
{2!\; 5!}=168\)}{\(\frac{8!}
{5!}=336\)}{\(8!=40\:320\)}{}{}}
\End

\Data\ID{9000139307}
\Updated{2021-04-25 06:04:56}
\Level{A}
\SubArea{Combinatorics}
\LANGen{
\Question{
There are \(10\) 
students hurrying to the dining room. Find the number of possibilities how can be
these students ordered into the row to get the meal.}
{\(10!=3\:628\:800\)}{\(10\)}{\(10^{10}\)}{\(\frac{10!}
{10} =362\:880\)}{}{}}
\LANGpl{
\Question{
\(10\) 
uczniów zmierza w kierunku stołówki. Wskaż ile jest możliwości ustawienia studentów w kolejce po posiłek.}
{\(10!=3\:628\:800\)}{\(10\)}{\(10^{10}\)}{\(\frac{10!}
{10}=362\:880 \)}{}{}}
\LANGcs{
\Question{
Na oběd běží \(10\) žáků. Kolika způsoby si mohou stoupnout do fronty k výdejnímu okénku?}
{\(10!=3\:628\:800\)}{\(10\)}{\(10^{10}\)}{\(\frac{10!}
{10} =362\:880\)}{}{}}
\LANGsk{
\Question{
Na obed beží \(10\) 
študentov. Koľkými spôsobmi sa môžu postaviť do rady k výdajnému okienku?}
{\(10!=3\:628\:800\)}{\(10\)}{\(10^{10}\)}{\(\frac{10!}
{10} =362\:880\)}{}{}}
\LANGes{
\Question{
\(10\) 
estudiantes están corriendo hacia el comedor del instituto. Halla el número de posibilidades distintas en las que se pueden ordenar
estos estudiantes en la fila para pedir la comida.}
{\(10!=3\:628\:800\)}{\(10\)}{\(10^{10}\)}{\(\frac{10!}
{10} =362\:880\)}{}{}}
\End

\Data\ID{9000139308}
\Updated{2021-04-25 07:04:36}
\Level{A}
\SubArea{Combinatorics}
\LANGen{
\Question{
The shooting club has \(25\) 
members. Among the members it is necessary to vote a board: a president, a cashier
and a webmaster. One person cannot have more than one of these positions and there
is only one member skilled enough to be a webmaster. How many possibilities exist to
set up the board?}
{\(24\cdot 23=552\)}{\(25\cdot 24=600\)}{\(24\cdot 23\cdot 22=12\:144\)}{\(25\cdot 24\cdot 23=13\:800\)}{}{}}
\LANGpl{
\Question{
Klub strzelców ma  \(25\) 
członków. Członkowie wybierają zarząd: prezesa, skarbnika
i sekretarza. Jedna osoba nie może mieć więcej niż jedną z tych funkcji, tylko jeden członek ma kwalifikacje, do pełnienia funkcji sekretarza. Ile jest możliwych sposobów wybrania zarządu?}
{\(24\cdot 23=552\)}{\(25\cdot 24=600\)}{\(24\cdot 23\cdot 22=12\:144\)}{\(25\cdot 24\cdot 23=13\:800\)}{}{}}
\LANGcs{
\Question{
Ve střeleckém klubu je \(25\) 
členů. Kolika způsoby z nich lze vybrat předsedu, pokladníka a správce www
stránek klubu, jestliže spravovat www stránky umí jen jeden z nich? Žádný
z členů nemůže zastávat více než jednu z uvedených funkcí.}
{\(24\cdot 23=552\)}{\(25\cdot 24=600\)}{\(24\cdot 23\cdot 22=12\:144\)}{\(25\cdot 24\cdot 23=13\:800\)}{}{}}
\LANGsk{
\Question{
Strelecký klub má \(25\) 
členov. Spomedzi členov je potrebné si zvoliť výbor zložený z prezidenta, pokladníka a správcu webových stránok. Žiadny z členov nemôže zastávať viac než jednu z uvedených funkcií a len jediný člen vie spravovať webové stránky. Koľkými spôsobmi je možné vytvoriť výbor?}
{\(24\cdot 23=552\)}{\(25\cdot 24=600\)}{\(24\cdot 23\cdot 22=12\:144\)}{\(25\cdot 24\cdot 23=13\:800\)}{}{}}
\LANGes{
\Question{
Un club de tiro tiene \(25\)
miembros. Entre los miembros es necesario votar una junta: un presidente, un secretario
y un webmaster. Una persona solo puede tener uno de estos puestos y solo un miembro está suficientemente capacitado para ser webmaster. ¿Cuántas posibilidades para la junta existen?}
{\(24\cdot 23=552\)}{\(25\cdot 24=600\)}{\(24\cdot 23\cdot 22=12\:144\)}{\(25\cdot 24\cdot 23=13\:800\)}{}{}}
\End

\Data\ID{9000139309}
\Updated{2021-04-25 07:04:58}
\Level{A}
\SubArea{Combinatorics}
\LANGen{
\Question{
There are \(20\) tablets in an
e-shop. From this amount \(18\) 
tablets are new and \(2\) 
tablets have been returned by customers. The e-shop manager gets an order
containing three tablets and he wants to get rid of the returned tablets first. How
many possibilities exist to complete the order?}
{\(18\)}{\(\frac{18!}
{3!\; 15!}=816\)}{\(18\cdot 16\cdot 3=864\)}{\(20\cdot 19\cdot 18=6\:840\)}{}{}}
\LANGpl{
\Question{
W sklepie elektronicznym znajduje się  \(20\) tabletów z  czego \(18\) 
tabletów jest nowych, a \(2\) 
zostały zwrócone przez klientów. Kierownik sklepu internetowego otrzymuje zamówienie
na  trzy tablety, w pierwszej kolejności chce pozbyć się zwróconych  tabletów. Ile jest  sposobów
realizacji tego  zamówienia?}
{\(18\)}{\(\frac{18!}
{3!\; 15!}=816\)}{\(18\cdot 16\cdot 3=864\)}{\(20\cdot 19\cdot 18=6\:840\)}{}{}}
\LANGcs{
\Question{
V e-shopu mají skladem \(20\) 
tabletů, z nichž \(18\) je
nových a \(2\) jsou vráceny
zákazníkem po \(14\) 
dnech používání. Zaměstnanec e-shopu má od majitele za úkol zbavit se
nejdříve použitých tabletů. Kolika způsoby může tento zaměstnanec
vybrat do objednávky nového zákazníka tři tablety tak, aby mezi nimi byly
oba použité a jeden nový?}
{\(18\)}{\(\frac{18!}
{3!\; 15!}=816\)}{\(18\cdot 16\cdot 3=864\)}{\(20\cdot 19\cdot 18=6\:840\)}{}{}}
\LANGsk{
\Question{
V e-shope majú \(20\) tabletov, z ktorých \(18\) je nových a \(2\) tablety sú vrátené zákazníkmi. Zamestnanec e-shopu má od majiteľa za úlohu zbaviť sa najprv použitých tabletov. Koľkými spôsobmi môže  vybrať do objednávky nového zákazníka tri tablety tak, aby medzi nimi boli obidva použité a jeden nový?}
{\(18\)}{\(\frac{18!}
{3!\; 15!}=816\)}{\(18\cdot 16\cdot 3=864\)}{\(20\cdot 19\cdot 18=6\:840\)}{}{}}
\LANGes{
\Question{
Hay \(20\) tabletas en una
tienda electrónica. De esta cantidad \(18\)
tabletas son nuevas y \(2\) tabletas han sido devueltas por los clientes.
 El gerente de la tienda electrónica recibe un pedido de tres tabletas y primero quiere usar las tabletas devueltas para este pedido. 
¿Cuántas posibilidades existen para organizar el pedido?}
{\(18\)}{\(\frac{18!}
{3!\; 15!}=816\)}{\(18\cdot 16\cdot 3=864\)}{\(20\cdot 19\cdot 18=6\:840\)}{}{}}
\End

\Data\ID{9000139310}
\Updated{2021-04-25 07:04:06}
\Level{A}
\SubArea{Combinatorics}
\LANGen{
\Question{
There are \(20\) tablets in an
e-shop. From this amount \(18\) 
tablets are new and \(2\) 
tablets have been returned by customers. The e-shop manager gets an order
containing three tablets and he wants to use only the new tablets for this order. How
many possibilities exist to complete the order?}
{\(\frac{18!}
{3!\; 15!}\)}{\(18\)}{\(18\cdot 16\cdot 3\)}{\(20\cdot 19\cdot 18\)}{}{}}
\LANGpl{
\Question{
W sklepie elektronicznym znajduje się 20 tabletów z czego 18 tabletów jest nowych, a 2 zostały zwrócone przez klientów. Kierownik sklepu internetowego otrzymuje zamówienie na trzy tablety, chce sprzedać tylko nowe tablety. Ile jest sposobów realizacji tego zamówienia?}
{\(\frac{18!}
{3!\; 15!}\)}{\(18\)}{\(18\cdot 16\cdot 3\)}{\(20\cdot 19\cdot 18\)}{}{}}
\LANGcs{
\Question{
V e-shopu mají skladem \(20\) 
tabletů, z nichž \(18\) je
nových a \(2\) jsou vráceny
zákazníkem po \(14\) 
dnech používání. Kolika způsoby může tento zaměstnanec vybrat do
objednávky nového zákazníka tři tablety tak, aby mezi nimi byly pouze nové
tablety?}
{\(\frac{18!}
{3!\; 15!}\)}{\(18\)}{\(18\cdot 16\cdot 3\)}{\(20\cdot 19\cdot 18\)}{}{}}
\LANGsk{
\Question{
V e-shope majú skladom \(20\) tabletov, z ktorých \(18\) 
 je nových a \(2\) 
tablety sú vrátené zákazníkmi. Koľkými spôsobmi môže zamestnanec vybrať do objednávky pre nového zákazníka tri tablety tak, aby medzi nimi boli len nové tablety?}
{\(\frac{18!}
{3!\; 15!}\)}{\(18\)}{\(18\cdot 16\cdot 3\)}{\(20\cdot 19\cdot 18\)}{}{}}
\LANGes{
\Question{
Hay \(20\)  tabletas en una tienda electrónica. De esta cantidad \(18\) 
tabletas son nuevas y \(2\) 
tabletas han sido devueltas por los clientes. El gerente de la tienda electrónica recibe un pedido de tres tabletas y primero quiere usar solo las nuevas tabletas para este pedido. ¿Cuántas posibilidades existen para organizar el pedido?}
{\(\frac{18!}
{3!\; 15!}\)}{\(18\)}{\(18\cdot 16\cdot 3\)}{\(20\cdot 19\cdot 18\)}{}{}}
\End

\Data\ID{9000139701}
\Updated{2021-04-25 07:04:59}
\Level{A}
\SubArea{Combinatorics}
\LANGen{
\Question{
There are \(15\) 
athletes in an athletic meeting. Determine in how many ways it is possible to
obtain the results on the first six places of the scoreboard if the place on
scoreboard cannot be shared (one athlete per one place on scoreboard).}
{\(\frac{15!}
 {9!} =3\:603\:600\)}{\(6^{15}=470\:184\:984\:576\)}{\(15!\, 6!=941\:525\:544\:960\:000\)}{\(\frac{15!}
{9!\, 6!}=5\:005\)}{}{}}
\LANGpl{
\Question{
W zawodach bierze udział \(15\) 
zawodników. Określ, ile jest możliwych sposobów
uzyskania   pierwszych sześciu miejsc na  tablicy wyników, jeśli miejsce na
tablicy wyników nie może być dzielone (jeden zawodnik na jednym miejscu na tablicy wyników).}
{\(\frac{15!}
 {9!}=3\:603\:600 \)}{\(6^{15}=470\:184\:984\:576\)}{\(15!\, 6!=941\:525\:544\:960\:000\)}{\(\frac{15!}
{9!\, 6!}=5\:005\)}{}{}}
\LANGcs{
\Question{
Soutěže se zúčastní \(15\) 
závodníků. Určete, kolika způsoby může být obsazeno prvních šest
bodovaných míst, pokud se na každém bodovaném místě umístí
právě jeden závodník.}
{\(\frac{15!}
 {9!} =3\:603\:600\)}{\(6^{15}=470\:184\:984\:576\)}{\(15!\, 6!=941\:525\:544\:960\:000\)}{\(\frac{15!}
{9!\, 6!}=5\:005\)}{}{}}
\LANGsk{
\Question{
Súťaže sa zúčastní \(15\) atlétov. Určte, koľkými spôsobmi môže byť obsadených prvých šesť bodovaných miest, ak sa na každom z nich umiestni práve jeden závodník.}
{\(\frac{15!}
 {9!} =3\:603\:600\)}{\(6^{15}=470\:184\:984\:576\)}{\(15!\, 6!=941\:525\:544\:960\:000\)}{\(\frac{15!}
{9!\, 6!}=5\:005\)}{}{}}
\LANGes{
\Question{
Hay \(15\)
atletas en una competición de atletismo. Determina de cuántas formas es posible
ocupar los primeros seis lugares de la clasificación si no es posible empatar.}
{\(\frac{15!}
 {9!} =3\:603\:600\)}{\(6^{15}=470\:184\:984\:576\)}{\(15!\, 6!=941\:525\:544\:960\:000\)}{\(\frac{15!}
{9!\, 6!}=5\:005\)}{}{}}
\End

\Data\ID{9000139702}
\Updated{2021-04-25 07:04:29}
\Level{A}
\SubArea{Combinatorics}
\LANGen{
\Question{
A race was attended by \(12\) 
competitors. Find the number of possible orderings of the competitors on the first
three places of the scoreboard.}
{\(\frac{12!}
 {9!} =1\:320\)}{\(3^{12}=531\:441\)}{\(\frac{12!}
{9!\, 3!}=220\)}{\(12!\, 3!=2\:874\:009\:600\)}{}{}}
\LANGpl{
\Question{
W wyścigu wzięło udział \(12 \)
konkurentów. Wskaż ilość  możliwości zdobycia przez nich  trzech pierwszych miejsc.}
{\({\frac{12!} {9!}}=1\:320\)}{\(3^{12}=531\:441\)}{\({\frac{12!}{9!\, 3!}}=220\)}{\(12!\, 3!=2\:874\:009\:600\)}{}{}}
\LANGcs{
\Question{
Závodů se zúčastnilo \(12\) 
závodníků. Určete, kolika způsoby mohou být uděleny zlatá, stříbrná a
bronzová medaile, pokud každou z medailí může získat pouze jeden
závodník.}
{\(\frac{12!}
 {9!} =1\:320\)}{\(3^{12}=531\:441\)}{\(\frac{12!}
{9!\, 3!}=220\)}{\(12!\, 3!=2\:874\:009\:600\)}{}{}}
\LANGsk{
\Question{
Pretekov sa zúčastnilo \(12\) pretekárov. Určte, koľkými spôsobmi sa môžu pretekári umiestniť na prvých troch miestach na výsledkovej tabuli.}
{\(\frac{12!}
 {9!}=1\:320 \)}{\(3^{12}=531\:441\)}{\(\frac{12!}
{9!\, 3!}=220\)}{\(12!\, 3!=2\:874\:009\:600\)}{}{}}
\LANGes{
\Question{
\(12\)
competidores asistieron a una carrera. ¿De cuántas maneras pueden los competidores recibir medalla de oro, de plata y de bronce?}
{\(\frac{12!}
 {9!} =1\:320\)}{\(3^{12}=531\:441\)}{\(\frac{12!}
{9!\, 3!}=220\)}{\(12!\, 3!=2\:874\:009\:600\)}{}{}}
\End

\Data\ID{9000139703}
\Updated{2021-04-25 07:04:02}
\Level{A}
\SubArea{Combinatorics}
\LANGen{
\Question{
The box contains \(5\) 
red crayons, \(4\) 
yellow crayons and \(2\) 
green crayons. The crayons are removed from the box and arranged in a
line. How many different color patterns can be obtained by this procedure?}
{\(\frac{11!}
{5!\, 4!\, 2!}=6\:930\)}{\(5\cdot 4\cdot 2=40\)}{\(5!\, 4!\, 2!=5\:760\)}{\(\left (5!\, 4!\right )^{2}=8\:294\:400\)}{}{}}
\LANGpl{
\Question{
W pudełku znajdują się kredki \(5\) 
czerwonych, \(4\) 
żółte i \(2\) 
zielone. Kredki zostały wyciągnięte z pudełka i ułożone w linię. Na ile możliwych sposobów można je ułożyć?}
{\(\frac{11!}
{5!\, 4!\, 2!}=6\:930\)}{\(5\cdot 4\cdot 2=40\)}{\(5!\, 4!\, 2!=5\:760\)}{\(\left (5!\, 4!\right )^{2}=8\:294\:400\)}{}{}}
\LANGcs{
\Question{
V krabičce je \(5\) 
červených pastelek, \(4\) 
žluté a \(2\) 
zelené pastelky. Určete, kolik různých barevných vzorů můžeme získat,
vyskládáme-li pastelky vedle sebe do krabičky.}
{\(\frac{11!}
{5!\, 4!\, 2!}=6\:930\)}{\(5\cdot 4\cdot 2=40\)}{\(5!\, 4!\, 2!=5\:760\)}{\(\left (5!\, 4!\right )^{2}=8\:294\:400\)}{}{}}
\LANGsk{
\Question{
Krabička obsahuje \(5\) červených, \(4\) žlté a \(2\) zelené pastelky. Pastelky vyberieme z krabičky a zoradíme ich vedľa seba. Koľko rôznych farebných vzorov môžeme takto získať?}
{\(\frac{11!}
{5!\, 4!\, 2!}=6\:930\)}{\(5\cdot 4\cdot 2=40\)}{\(5!\, 4!\, 2!=5\:760\)}{\(\left (5!\, 4!\right )^{2}=8\:294\:400\)}{}{}}
\LANGes{
\Question{
Una caja contiene \(5\)
lápices de color rojo, \(4\)
lápices de color amarillo y \(2\)
lápices de color verde. Los lápices de colores se sacan de la caja y se colocan en una
línea. ¿Cuántos patrones diferentes de color se pueden obtener con este procedimiento?}
{\(\frac{11!}
{5!\, 4!\, 2!}=6\:930\)}{\(5\cdot 4\cdot 2=40\)}{\(5!\, 4!\, 2!=5\:760\)}{\(\left (5!\, 4!\right )^{2}=8\:294\:400\)}{}{}}
\End

\Data\ID{9000139704}
\Updated{2021-04-25 07:04:49}
\Level{A}
\SubArea{Combinatorics}
\LANGen{
\Question{
There are \(5\) 
different kinds of cakes in a shop. Find the number of possibilities how to buy
\(8\) cakes in this shop.
(There is more than \(8\) 
cakes of each kind available.)}
{\(\frac{12!}
{8!\, 4!}=495\)}{\(5!\, 8!=4\:838\:400\)}{\(5^{8}=390\:625\)}{\(\frac{8!}
{5!\, 3!}=56\)}{}{}}
\LANGpl{
\Question{
W cukierni jest  \(5\) 
różnych rodzajów tortów. Wskaż ile jest możliwości kupienia 
\(8\) tortów
(Jest więcej niż \(8\) 
tortów z każdego rodzaju.)}
{\(\frac{12!}
{8!\, 4!}=495\)}{\(5!\, 8!=4\:838\:400\)}{\(5^{8}=390\:625\)}{\(\frac{8!}
{5!\, 3!}=56\)}{}{}}
\LANGcs{
\Question{
V cukrárně je \(5\) 
druhů zákusků v dostatečně velkém množství. Určete, kolika způsoby můžeme
koupit \(8\) 
zákusků.}
{\(\frac{12!}
{8!\, 4!}=495\)}{\(5!\, 8!=4\:838\:400\)}{\(5^{8}=390\:625\)}{\(\frac{8!}
{5!\, 3!}=56\)}{}{}}
\LANGsk{
\Question{
V cukrárni je \(5\) 
rôznych druhov zákuskov v dostatočne veľkom množstve. Určte, koľkými spôsobmi môžeme kúpiť 
\(8\) zákuskov.}
{\(\frac{12!}
{8!\, 4!}=495\)}{\(5!\, 8!=4\:838\:400\)}{\(5^{8}=390\:625\)}{\(\frac{8!}
{5!\, 3!}=56\)}{}{}}
\LANGes{
\Question{
Hay \(5\)
diferentes tipos de pasteles en una tienda. Halla el número de posibilidades para comprar
\(8\) pasteles en esta tienda.
(Hay más de \(8\)
pasteles de cada tipo disponibles.)}
{\(\frac{12!}
{8!\, 4!}=495\)}{\(5!\, 8!=4\:838\:400\)}{\(5^{8}=390\:625\)}{\(\frac{8!}
{5!\, 3!}=56\)}{}{}}
\End

\Data\ID{9000139705}
\Updated{2021-04-25 07:04:39}
\Level{A}
\SubArea{Combinatorics}
\LANGen{
\Question{
From the group of \(10\) boys
and \(5\) girls we have to
select a small group of \(3\) 
boys and \(2\) 
girls. How many possibilities exist for this choice?}
{\(\frac{10!}
{7!\, 3!}\cdot \frac{5!}
{3!\, 2!}=1\:200\)}{\(5^{10}=9\:765\:625\)}{\(10\cdot 5!\, 3!=7\:200\)}{\(5\cdot \frac{10!}
 {3!} =3\:024\:000\)}{}{}}
\LANGpl{
\Question{
Z grupy  \(10\) chłopców
i \(5\) dziewczynek musimy wybrać grupę  \(3\) 
chłopców i \(2\) 
dziewczynek. Ile istnieje możliwości dokonania tego wyboru?}
{\(\frac{10!}
{7!\, 3!}\cdot \frac{5!}
{3!\, 2!}=1\:200\)}{\(5^{10}=9\:765\:625\)}{\(10\cdot 5!\, 3!=7\:200\)}{\(5\cdot \frac{10!}
 {3!} =3\:024\:000\)}{}{}}
\LANGcs{
\Question{
Určete, kolika způsoby můžeme z
\(10\) chlapců a
\(5\) děvčat vybrat
pětici, ve které budou \(3\) 
chlapci a dvě děvčata.}
{\(\frac{10!}
{7!\, 3!}\cdot \frac{5!}
{3!\, 2!}=1\:200\)}{\(5^{10}=9\:765\:625\)}{\(10\cdot 5!\, 3!=7\:200\)}{\(5\cdot \frac{10!}
 {3!} =3\:024\:000\)}{}{}}
\LANGsk{
\Question{
Koľkými spôsobmi môžeme zo skupiny \(10\) chlapcov a \(5\) dievčat vybrať päťčlenné družstvo, v ktorom budú \(3\) chlapci a \(2\) dievčatá?}
{\(\frac{10!}
{7!\, 3!}\cdot \frac{5!}
{3!\, 2!}=1\:200\)}{\(5^{10}=9\:765\:625\)}{\(10\cdot 5!\, 3!=7\:200\)}{\(5\cdot \frac{10!}
 {3!} =3\:024\:000\)}{}{}}
\LANGes{
\Question{
De un grupo de \(10\) chicos
y \(5\) chicas tenemos que
seleccionar un  subgrupo de \(3\) 
chicos y \(2\) 
chicas.¿Cuántas posibilidades existen para esta selección?}
{\(\frac{10!}
{7!\, 3!}\cdot \frac{5!}
{3!\, 2!}=1\:200\)}{\(5^{10}=9\:765\:625\)}{\(10\cdot 5!\, 3!=7\:200\)}{\(5\cdot \frac{10!}
 {3!} =3\:024\:000\)}{}{}}
\End

\Data\ID{9000139706}
\Updated{2021-04-25 07:04:17}
\Level{A}
\SubArea{Combinatorics}
\LANGen{
\Question{
The international alphabet contains \(26\) 
letters. The letters of this alphabet and the digits from
\(0\) to
\(9\) are used to form a
code of the length \(4\) 
(a code contains \(4\) 
characters). The characters may repeat through the code and the code is not case
sensitive (uppercase letters are equivalent to lowercase letters). How many codes can
be obtained?}
{\(36^{4}=1\:679\:616\)}{\(10\cdot 26^{4}=4\:569\:760\)}{\(\frac{36!}
{32!\, 4!}=58\:905\)}{\(\frac{26!}
{22!\, 4!}=14\:950\)}{}{}}
\LANGpl{
\Question{
Międzynarodowy alfabet zawiera 26
liter. Litery tego alfabetu i cyfry od
0  do
9  zostały użyte  do utworzenia 
kodu o  długości 4 
(kod zawiera 4
znaki). Znaki mogą powtarzać się, kod nie uwzględnia wielkość liter. Ile kodów można
uzyskać?}
{\(36^{4}=1\:679\:616\)}{\(10\cdot 26^{4}=4\:569\:760\)}{\(\frac{36!}
{32!\, 4!}=58\:905\)}{\(\frac{26!}
{22!\, 4!}=14\:950\)}{}{}}
\LANGcs{
\Question{
Mezinárodní abeceda má \(26\) 
písmen. Určete počet možností čtyřmístného kódu tvořeného
malými písmeny této abecedy a číslicemi od \(0\) do
\(9\).  Znaky se mohou opakovat.}
{\(36^{4}=1\:679\:616\)}{\(10\cdot 26^{4}=4\:569\:760\)}{\(\frac{36!}
{32!\, 4!}=58\:905\)}{\(\frac{26!}
{22!\, 4!}=14\:950\)}{}{}}
\LANGsk{
\Question{
Medzinárodná abeceda má \(26\) 
písmen. Určte počet možností štvormiestneho kódu tvoreného malými písmenami tejto abecedy a číslicami od
\(0\) do
\(9\). Znaky sa môžu opakovať.}
{\(36^{4}=1\:679\:616\)}{\(10\cdot 26^{4}=4\:569\:760\)}{\(\frac{36!}
{32!\, 4!}=58\:905\)}{\(\frac{26!}
{22!\, 4!}=14\:950\)}{}{}}
\LANGes{
\Question{
El alfabeto internacional tiene \(26\)
letras. Calcula el número de opciones para un código de cuatro dígitos formado por las
letras minúsculas de este alfabeto y números de \(0\) a
\(9\). Los caracteres se pueden repetir.}
{\(36^{4}=1\:679\:616\)}{\(10\cdot 26^{4}=4\:569\:760\)}{\(\frac{36!}
{32!\, 4!}=58\:905\)}{\(\frac{26!}
{22!\, 4!}=14\:950\)}{}{}}
\End

\Data\ID{9000139707}
\Updated{2021-04-25 07:04:40}
\Level{A}
\SubArea{Combinatorics}
\LANGen{
\Question{
A Morse code utilized dots and dashes to encode letters of
an alphabet. Find the number of signals of the length from
\(1\) to
\(4\) which
can be obtained from dots and dashes.}
{\(2 + 2^{2} + 2^{3} + 2^{4}=30\)}{\(1 + 2 + 3! + 4!=33\)}{\(\frac{4!}
{3!\, 2!}=2\)}{\(2 \cdot 1 + 2 \cdot 2 + 2 \cdot 3 + 2 \cdot 4=20\)}{}{}}
\LANGpl{
\Question{
Alfabet Morse'a używa kropek i myślników. Określ liczbę  sygnałów o długości od jeden do cztery, które mogą zostać utworzone z kropek i myślników.}
{\(2 + 2^{2} + 2^{3} + 2^{4}=30\)}{\(1 + 2 + 3! + 4!=33\)}{\(\frac{4!}
{3!\, 2!}=2\)}{\(2 \cdot 1 + 2 \cdot 2 + 2 \cdot 3 + 2 \cdot 4=20\)}{}{}}
\LANGcs{
\Question{
Morseova abeceda používá tečky a čárky. Určete počet jednomístných
až čtyřmístných skupin tvořených pomocí teček a čárek.}
{\(2 + 2^{2} + 2^{3} + 2^{4}=30\)}{\(1 + 2 + 3! + 4!=33\)}{\(\frac{4!}
{3!\, 2!}=2\)}{\(2 \cdot 1 + 2 \cdot 2 + 2 \cdot 3 + 2 \cdot 4=20\)}{}{}}
\LANGsk{
\Question{
Morseova abeceda používa znaky, ktoré sú vytvorené pomocou bodiek a čiarok. Určte, koľko je v nej jednomiestnych až štvormiestnych znakov.}
{\(2 + 2^{2} + 2^{3} + 2^{4}=30\)}{\(1 + 2 + 3! + 4!=33\)}{\(\frac{4!}
{3!\, 2!}=2\)}{\(2 \cdot 1 + 2 \cdot 2 + 2 \cdot 3 + 2 \cdot 4=20\)}{}{}}
\LANGes{
\Question{
Un código Morse utiliza puntos y rayas para codificar letras de
un alfabeto. Halla el número de señales de longitud de
\(1\) a
\(4\) que
se pueden obtener mediante puntos y rayas.}
{\(2 + 2^{2} + 2^{3} + 2^{4}=30\)}{\(1 + 2 + 3! + 4!=33\)}{\(\frac{4!}
{3!\, 2!}=2\)}{\(2 \cdot 1 + 2 \cdot 2 + 2 \cdot 3 + 2 \cdot 4=20\)}{}{}}
\End

\Data\ID{9000139708}
\Updated{2021-04-25 07:04:50}
\Level{A}
\SubArea{Combinatorics}
\LANGen{
\Question{
The shelf contains \(15\) 
books. From this amount, \(9\) 
books are in English and \(6\) 
books in other languages. Find the number of possibilities how to rearrange the
books on the shelf, if all English books have to be on the left and the other on the
right.}
{\(9!\, 6!=261\:273\:600\)}{\(9^{6}=531\:441\)}{\(\frac{9!}
{6!}=504\)}{\(\frac{9!}
{6!\, 3!}=84\)}{}{}}
\LANGpl{
\Question{
Na półce znajduje się \(15\) 
książek, z czego  \(9\) 
książek jest w języku angielskim, \(6\) 
w innych językach. Ile jest możliwości ułożenia książek, jeżeli wszystkie książki w języku angielskim muszą być ustawione po lewej stronie, a pozostałe po prawej?}
{\(9!\, 6!=261\:273\:600\)}{\(9^{6}=531\:441\)}{\(\frac{9!}
{6!}=504\)}{\(\frac{9!}
{6!\, 3!}=84\)}{}{}}
\LANGcs{
\Question{
Na polici je \(9\) různých
knih v češtině a \(6\) 
různých knih cizojazyčných. Určete, kolika způsoby můžeme knihy
přeskládat tak, aby za sebou byly seřazeny nejprve česky psané knihy a za
nimi knihy cizojazyčné.}
{\(9!\, 6!=261\:273\:600\)}{\(9^{6}=531\:441\)}{\(\frac{9!}
{6!}=504\)}{\(\frac{9!}
{6!\, 3!}=84\)}{}{}}
\LANGsk{
\Question{
Na polici  je \(15\) 
kníh, z toho je \(9\) rôznych kníh
v slovenčine a \(6\) rôznych kníh cudzojazyčných. 
Určte, koľkými spôsobmi môžeme knihy usporiadať tak, aby za sebou boli zoradené najprv slovenské a za nimi cudzojazyčné knihy.}
{\(9!\, 6!=261\:273\:600\)}{\(9^{6}=531\:441\)}{\(\frac{9!}
{6!}=504\)}{\(\frac{9!}
{6!\, 3!}=84\)}{}{}}
\LANGes{
\Question{
En un estante se encuentran  \(15\) 
libros. De esta cantidad, \(9\) 
 libros están escritos en inglés y \(6\) 
libros en otros idiomas. Halla  el número de posibilidades para reorganizar los
libros en el estante si todos los libros escritos en inglés tienen que estar a la izquierda y los otros a la derecha.}
{\(9!\, 6!=261\:273\:600\)}{\(9^{6}=531\:441\)}{\(\frac{9!}
{6!}=504\)}{\(\frac{9!}
{6!\, 3!}=84\)}{}{}}
\End

\Data\ID{9000139709}
\Updated{2021-04-25 07:04:28}
\Level{A}
\SubArea{Combinatorics}
\LANGen{
\Question{
Find the number of possible ways how to organize a group of
\(10\) 
children in one ordered row if Mike and Paul should be next to each other.}
{\(2\cdot 9!=725\:760\)}{\(10!=3\:628\:800\)}{\(\frac{10!}
 {2!} =1\:814\:400\)}{\(\frac{10!}
 {8!} =90\)}{}{}}
\LANGpl{
\Question{
Ile jest sposobów ustawienia 
\(10\) 
dzieci w jednym rzędzie,  jeżeli Adam i Alek muszą stać obok siebie?}
{\(2\cdot 9!=725\:760\)}{\(10!=3\:628\:800\)}{\(\frac{10!}
 {2!} =1\:814\:400\)}{\(\frac{10!}
 {8!} =90\)}{}{}}
\LANGcs{
\Question{
Určete, kolika způsoby můžeme postavit
\(10\) 
dětí vedle sebe tak, aby Adam a Bedřich stáli vedle sebe.}
{\(2\cdot 9!=725\:760\)}{\(10!=3\:628\:800\)}{\(\frac{10!}
 {2!}=1\:814\:400 \)}{\(\frac{10!}
 {8!} =90\)}{}{}}
\LANGsk{
\Question{
Určte, koľkými spôsobmi môžeme postaviť
\(10\) detí vedľa seba tak, aby Michal a Pavol stáli vedľa seba.}
{\(2\cdot 9!=725\:760\)}{\(10!=3\:628\:800\)}{\(\frac{10!}
 {2!}=1\:814\:400 \)}{\(\frac{10!}
 {8!}=90 \)}{}{}}
\LANGes{
\Question{
Halla el número de posibilidades para organizar un grupo de
\(10\)
niños en una fila para que  Mike y Paul estén uno al lado del otro.}
{\(2\cdot 9!=725\:760\)}{\(10!=3\:628\:800\)}{\(\frac{10!}
 {2!} =1\:814\:400\)}{\(\frac{10!}
 {8!} =90\)}{}{}}
\End

\Data\ID{9000148901}
\Updated{2021-04-25 07:04:55}
\Level{A}
\SubArea{Combinatorics}
\LANGen{
\Question{
The current Czech vehicle registration plate number has the form NLN-NNNN, where N stands
for a digit from \(0\) 
to \(9\) 
and L stands for a letter from an alphabet containing
\(26\) 
letters. How many different registration plates are possible?}
{\(26\cdot 10^{6}\)}{\(10^{6}\)}{\(15\cdot 10^{6} + 6\cdot 10^{5}= 156\cdot 10^{5}\)}{\(16\cdot 10^{6}\)}{}{}}
\LANGpl{
\Question{
Czeska tablica rejestracyjna pojazdów ma postać NLN-NNNN, gdzie N oznacza
cyfrę od \(0\) 
do \(9\) 
a L oznacza literę alfabetu zawierającego 
\(26\) liter. Ile jest możliwości utworzenia tablicy rejestracyjnej?}
{\(26\cdot 10^{6}\)}{\(10^{6}\)}{\(15\cdot 10^{6} + 6\cdot 10^{5}= 156\cdot 10^{5}\)}{\(16\cdot 10^{6}\)}{}{}}
\LANGcs{
\Question{
V současnosti používané státní poznávací značky
automobilů mají tvar CPC-CCCC, kde C označuje číslici od
\(0\) do
\(9\) a P písmeno z
mezinárodní abecedy s \(26\) 
znaky. Kolik státních poznávacích značek v uvedeném tvaru je možné
sestavit?}
{\(26\cdot 10^{6}\)}{\(10^{6}\)}{\(15\cdot 10^{6} + 6\cdot 10^{5}= 156\cdot 10^{5}\)}{\(16\cdot 10^{6}\)}{}{}}
\LANGsk{
\Question{
V súčasnosti používané štátne poznávacie značky českých automobilov majú tvar CPC-CCCC, kde C označuje 
číslicu od \(0\) 
do \(9\) 
a P označuje písmeno abecedy obsahujúcej 
\(26\) 
písmen. Koľko rôznych štátnych poznávacích značiek v uvedenom tvare je možné zostaviť?}
{\(26\cdot 10^{6}\)}{\(10^{6}\)}{\(15\cdot 10^{6} + 6\cdot 10^{5}= 156\cdot 10^{5}\)}{\(16\cdot 10^{6}\)}{}{}}
\LANGes{
\Question{
Las matrículas checas actuales de vehículos tienen la forma NLN-NNNN, donde N puede ser
un dígito de \(0\)
a \(9\)
y L representa una letra del alfabeto que contiene
\(26\)
letras. ¿Cuántas matrículas diferentes se pueden crear?}
{\(26\cdot 10^{6}\)}{\(10^{6}\)}{\(15\cdot 10^{6} + 6\cdot 10^{5}= 156\cdot 10^{5}\)}{\(16\cdot 10^{6}\)}{}{}}
\End

\Data\ID{9000148902}
\Updated{2020-08-19 12:08:40}
\Level{A}
\SubArea{Combinatorics}
\LANGen{
\Question{
There are four paths from a city to the top of nearby mountain. Find the number of
possible treks from the city to the mountain and back, if it is required to use one
path up and another one down.}
{\(4\cdot 3=12\)}{\(4\cdot 4=16\)}{\(4 + 3=7\)}{\(2\cdot 4=8\)}{}{}}
\LANGpl{
\Question{
Z miasta na szczyt pobliskiej góry prowadzą cztery ścieżki. Wskaż liczbę możliwych wycieczek z miasta na górę i z góry do miasta, jeżeli nie można wracać tą samą ścieżką.}
{\(4\cdot 3=12\)}{\(4\cdot 4=16\)}{\(4 + 3=7\)}{\(2\cdot 4=8\)}{}{}}
\LANGcs{
\Question{
Z Pece pod Sněžkou vedou na vrchol Sněžky
(\(1\: 602\, \mathrm{m}\)) v
podstatě čtyři cesty: lanovkou, přes Růžohorky, Obřím dolem a přes
Výrovku. Určete počet způsobů, kterými je možno se dostat na
vrchol a zpět tak, aby zpáteční cesta byla jiná než cesta na vrchol.}
{\(4\cdot 3=12\)}{\(4\cdot 4=16\)}{\(4 + 3=7\)}{\(2\cdot 4=8\)}{}{}}
\LANGsk{
\Question{
Na vrchol blízkej hory vedú z mesta štyri cesty. Určte počet spôsobov, ktorými je možné sa dostať na vrchol a späť tak, aby spiatočná cesta bola iná ako cesta na vrchol.}
{\(4\cdot 3=12\)}{\(4\cdot 4=16\)}{\(4 + 3=7\)}{\(2\cdot 4=8\)}{}{}}
\LANGes{
\Question{
Hay cuatro caminos desde una ciudad hasta la cima de una montaña cercana. Halla el número de
rutas posibles  desde la ciudad a la montaña y viceversa suponiendo que es necesario utilizar un
camino hacia arriba y otro hacia abajo.}
{\(4\cdot 3=12\)}{\(4\cdot 4=16\)}{\(4 + 3=7\)}{\(2\cdot 4=8\)}{}{}}
\End

\Data\ID{9000148903}
\Updated{2021-04-25 07:04:54}
\Level{A}
\SubArea{Combinatorics}
\LANGen{
\Question{
A combination lock will open if a right choice of three numbers (from
\(1\) to
\(9\)) is selected.
Suppose that we use a brute force attack to open the lock (we try all possibilities). To try one
code takes \(20\) 
seconds. What is the maximal time (in seconds) required to open the lock by brute
force?}
{\(20\cdot 9^{3}\, \mathrm{s}=14\:580\,\mathrm{s}\)}{\(20\cdot \frac{9!}
{6!}\, \mathrm{s}=10\:080\,\mathrm{s}\)}{\(20\cdot \frac{9!}
{3!\; 6!}\, \mathrm{s}=1\:680\,\mathrm{s}\)}{\(20\cdot 9\cdot 3\, \mathrm{s}=540\,\mathrm{s}\)}{}{}}
\LANGpl{
\Question{
Zamek zostanie otworzony jeżeli zostaną wybrane właściwe trzy cyfry (od
\(1\) do \(9\)).
Załóżmy, że używamy metody  "na siłę", aby otworzyć blokadę (spróbujemy wszystkich możliwości). Sprawdzenie jednego kodu trwa \(20\) 
sekund. Jaki jest maksymalny czas (w sekundach) potrzebny do otworzenia zamka metodą "na siłę"?}
{\(20\cdot 9^{3}\, \mathrm{s}=14\:580\,\mathrm{s}\)}{\(20\cdot \frac{9!}
{6!}\, \mathrm{s}=10\:080\,\mathrm{s}\)}{\(20\cdot \frac{9!}
{3!\; 6!}\, \mathrm{s}=1\:680\,\mathrm{s}\)}{\(20\cdot 9\cdot 3\, \mathrm{s}=540\,\mathrm{s}\)}{}{}}
\LANGcs{
\Question{
Kód zámku na kolo je trojmístný a skládá se z číslic od
\(1\) do
\(9\). Jak
dlouho budu odemykat zámek, když zapomenu kód a uhodnu kód až
posledním možným pokusem? Vytočení jednoho kódu trvá dvacet vteřin.}
{\(20\cdot 9^{3}\, \mathrm{s}=14\:580\,\mathrm{s}\)}{\(20\cdot \frac{9!}
{6!}\, \mathrm{s}=10\:080\,\mathrm{s}\)}{\(20\cdot \frac{9!}
{3!\; 6!}\, \mathrm{s}=1\:680\,\mathrm{s}\)}{\(20\cdot 9\cdot 3\, \mathrm{s}=540\,\mathrm{s}\)}{}{}}
\LANGsk{
\Question{
Kód zámku je trojmiestne číslo a skladá sa z číslic od
\(1\) do
\(9\). Ako dlho budeme odomykať zámok, keď zabudneme kód a uhádneme ho až posledným možným pokusom?
Vytočenie jedného kódu trvá  \(20\) 
sekúnd.}
{\(20\cdot 9^{3}\, \mathrm{s}=14\:580\,\mathrm{s}\)}{\(20\cdot \frac{9!}
{6!}\, \mathrm{s}=10\:080\,\mathrm{s}\)}{\(20\cdot \frac{9!}
{3!\; 6!}\, \mathrm{s}=1\:680\,\mathrm{s}\)}{\(20\cdot 9\cdot 3\, \mathrm{s}=540\,\mathrm{s}\)}{}{}}
\LANGes{
\Question{
Un candado de combinacón se abre  si se eligen correctamente tres números de
\(1\) a
\(9\).
¿Cuánto tiempo necesitas para desbloquear el candado si has olvidado el código y lo vas a adivinar en el último intento posible? Se necesitan \(20\) segundos para marcar un código.}
{\(20\cdot 9^{3}\, \mathrm{s}=14\:580\,\mathrm{s}\)}{\(20\cdot \frac{9!}
{6!}\, \mathrm{s}=10\:080\,\mathrm{s}\)}{\(20\cdot \frac{9!}
{3!\; 6!}\, \mathrm{s}=1\:680\,\mathrm{s}\)}{\(20\cdot 9\cdot 3\, \mathrm{s}=540\,\mathrm{s}\)}{}{}}
\End

\Data\ID{9000148904}
\Updated{2020-08-19 12:08:57}
\Level{A}
\SubArea{Combinatorics}
\LANGen{
\Question{
Pamela needs new ski for a ski course. There are skis from six different vendors in a
shop. The shop has four different ski pairs from each vendor, but two vendors have all
products behind Pam's financial limit. How many pairs are at disposal for Pam?}
{\(4\cdot 4=16\)}{\(4!=24\)}{\(4\cdot 2=8\)}{\(4 + 2=6\)}{}{}}
\LANGpl{
\Question{
Basia  potrzebuje nowych nart na kurs narciarski. W sklepie znajdują się narty od sześciu różnych producentów. W asortymencie sklepu znajdują się cztery różne  pary nart od każdego producenta, jednak narty od dwóch producentów są poza możliwościami finansowymi Basi. Ile par nart może kupić Basia?}
{\(4\cdot 4=16\)}{\(4!=24\)}{\(4\cdot 2=8\)}{\(4 + 2=6\)}{}{}}
\LANGcs{
\Question{
Veronika potřebuje na lyžařský kurz nové lyže. V obchodě mají \(6\) různých značek lyží a od každé z nich mají \(4\) různé páry. Z kolika párů lyží může Veronika vybírat, jsou-li všechny páry lyží dvou značek nad její finanční možnosti?}
{\(4\cdot 4=16\)}{\(4!=24\)}{\(4\cdot 2=8\)}{\(4 + 2=6\)}{}{}}
\LANGsk{
\Question{
Veronika potrebuje nové lyže na lyžiarsky kurz. V obchode majú 6 rôznych značiek lyží a od každej z nich majú štyri rôzne páry. Z koľkých párov lyží si môže Veronika vyberať, ak sú všetky páry lyží dvoch značiek nad jej finančné možnosti?}
{\(4\cdot 4=16\)}{\(4!=24\)}{\(4\cdot 2=8\)}{\(4 + 2=6\)}{}{}}
\LANGes{
\Question{
Pamela necesita esquís nuevos para un curso de esquí. Hay esquís de seis proveedores diferentes en una
tienda. La tienda tiene cuatro pares de esquís diferentes de cada proveedor, pero dos proveedores tienen todos
productos por encima de las posibilidades económicas de Pam. ¿Cuántos pares hay disponibles para Pam?}
{\(4\cdot 4=16\)}{\(4!=24\)}{\(4\cdot 2=8\)}{\(4 + 2=6\)}{}{}}
\End

\Data\ID{9000148905}
\Updated{2021-04-25 07:04:28}
\Level{A}
\SubArea{Combinatorics}
\LANGen{
\Question{
The king has eight daughters. A dragon asks two daughters, otherwise he will
destroy the whole country. In how many ways can we choose a pair of king's
daughters for the dragon?}
{\(\frac{8!}
{2!\; 6!}=28\)}{\(8\cdot 7=56\)}{\(8^{2}=64\)}{\(8 + 7=15\)}{}{}}
\LANGpl{
\Question{
Król ma osiem córek. Smok żąda dwóch córek w przeciwnym razie zniszczy całe miasto. Ile jest sposobów wybrania dwóch królewskich córek dla smoka?}
{\(\frac{8!}
{2!\; 6!}=28\)}{\(8\cdot 7=56\)}{\(8^{2}=64\)}{\(8 + 7=15\)}{}{}}
\LANGcs{
\Question{
Král má osm dcer. Určete, kolika způsoby může vybrat dvě dcery, které
chce sníst stohlavý drak. Není důležité, kterou princeznu vybereme jako
první a kterou jako druhou, protože drak bude jíst obě princezny najednou.}
{\(\frac{8!}
{2!\; 6!}=28\)}{\(8\cdot 7=56\)}{\(8^{2}=64\)}{\(8 + 7=15\)}{}{}}
\LANGsk{
\Question{
Kráľ má osem dcér. Drak žiada dve kráľove dcéry, inak zničí celú krajinu. Koľkými spôsobmi môžeme vybrať dve dcéry pre draka?}
{\(\frac{8!}
{2!\; 6!}=28\)}{\(8\cdot 7=56\)}{\(8^{2}=64\)}{\(8 + 7=15\)}{}{}}
\LANGes{
\Question{
Un rey tiene ocho hijas. Un dragón le pide dos hijas o destruirá todo el reino. ¿De cuántas formas podemos elegir un par de hijas del rey para el dragón?}
{\(\frac{8!}
{2!\; 6!}=28\)}{\(8\cdot 7=56\)}{\(8^{2}=64\)}{\(8 + 7=15\)}{}{}}
\End

\Data\ID{9000148907}
\Updated{2021-04-25 07:04:07}
\Level{A}
\SubArea{Combinatorics}
\LANGen{
\Question{
A bowl contains \(12\) different
gummy-bears and \(20\) 
different sweet-drops. Anne can choose either one sweet-drop or one gummy-bear.
From the rest, Jane can choose one sweet-drop and two gummy-bears. Anne wants to
provide a maximum of the possibilities for Jane's choice. What should Anne choose?}
{sweet-drop}{gummy-bear}{Both possibilities give the same result.}{}{}{}}
\LANGpl{
\Question{
W misce znajduje się  \(12\) różnych 
żelków i  \(20\) 
różnych dropsów. Ania może wybrać dropsa albo żelka.
Pozostali, Jan może wybrać jednego dropsa i dwa żelki. Ania chce, aby Jan miał maksymalną liczbę możliwości wyboru.
Co powinna wybrać Ania?}
{dropsa}{żelka}{Obie odpowiedzi mają ten sam wynik.}{}{}{}}
\LANGcs{
\Question{
V misce je \(12\) gumových
bonbonů a \(20\) 
hašlerek. Anička si může vybrat buď jednu hašlerku, anebo jeden gumový
bonbon tak, aby Pavla, která si po ní vybere jednu hašlerku a dva gumové
bonbony, měla co nejvíce různých možností výběru. Co si má Anička
vybrat?}
{Anička si musí vybrat hašlerku.}{Anička si musí vybrat gumový bonbon.}{Je to jedno, Anička si může vybrat, jakou sladkost chce.}{}{}{}}
\LANGsk{
\Question{
V miske je \(12\) rozdielnych gumových cukríkov a \(20\) hašleriek. Anička si môže vybrať jednu hašlerku alebo jeden gumový cukrík.
Zo zvyšku si Janka môže vybrať jednu hašlerku a dva gumové cukríky. Anička chce, aby Janka mala čo najviac rôznych možností výberu. Čo by si mala Anička vybrať?}
{Anička si musí vybrať hašlerku.}{Anička si musí vybrať gumový cukrík.}{Je to jedno, Anička si môže vybrať ľubovoľnú jednu sladkosť.}{}{}{}}
\LANGes{
\Question{
Un cuenco contiene \(12\) 
ositos de goma diferentes y \(20\) caramelos diferentes.
 Ana puede elegir entre un caramelo o un osito de goma.
Del resto, Jana puede elegir un caramelo y dos ositos de goma. Ana quiere
facilitar al máximo las posibilidades para la elección de Jane. ¿Qué debería elegir Ana?}
{caramelo}{osito de goma}{Ambas posibilidades daran el mismo resultado.}{}{}{}}
\End

\Data\ID{9000148908}
\Updated{2021-04-25 08:04:11}
\Level{A}
\SubArea{Combinatorics}
\LANGen{
\Question{
There are seven different yellow apples, eight different green apples and ten different
red apples. How many ways are there to choose three apples, if we wish to have three
apples of different colors?}
{\(10\cdot 8\cdot 7=560\)}{\(\frac{10\cdot 8\cdot 7}
 {2}=280 \)}{\((10 + 8 + 7)\cdot 2=50\)}{\(10 + 8 + 7=25\)}{}{}}
\LANGpl{
\Question{
Jest siedem różnych żółtych jabłek, osiem różnych zielonych jabłek i dziesięć różnych
czerwonych jabłek. Ile jest sposobów wyboru trzech jabłek, jeśli chcemy mieć trzy
jabłka o różnych kolorach?}
{\(10\cdot 8\cdot 7=560\)}{\(\frac{10\cdot 8\cdot 7}
 {2} =280\)}{\((10 + 8 + 7)\cdot 2=50\)}{\(10 + 8 + 7=25\)}{}{}}
\LANGcs{
\Question{
V misce je sedm různých žlutých jablek, osm různých zelených
jablek a deset různých červených jablek. Kolika způsoby lze provést
výběr tří jablek, jestliže chceme, aby každé jablko bylo jiné barvy?}
{\(10\cdot 8\cdot 7=560\)}{\(\frac{10\cdot 8\cdot 7}
 {2} =280\)}{\((10 + 8 + 7)\cdot 2=50\)}{\(10 + 8 + 7=25\)}{}{}}
\LANGsk{
\Question{
V miske je sedem rôznych žltých jabĺk, osem rôznych zelených jabĺk a desať  rôznych červených jabĺk. Koľkými spôsobmi môžeme vybrať tri jablká, ak chceme, aby každé jablko bolo inej farby?}
{\(10\cdot 8\cdot 7=560\)}{\(\frac{10\cdot 8\cdot 7}
 {2} =280\)}{\((10 + 8 + 7)\cdot 2=50\)}{\(10 + 8 + 7=25\)}{}{}}
\LANGes{
\Question{
Hay siete manzanas diferentes amarillas, ocho manzanas diferentes verdes y diez manzanas diferentes rojas. ¿De cuántas maneras hay que elegir tres manzanas, si deseamos tener tres
manzanas de diferentes colores?}
{\(10\cdot 8\cdot 7=560\)}{\(\frac{10\cdot 8\cdot 7}
 {2}=280 \)}{\((10 + 8 + 7)\cdot 2=50\)}{\(10 + 8 + 7=25\)}{}{}}
\End

\Data\ID{9000148909}
\Updated{2020-08-19 01:08:30}
\Level{A}
\SubArea{Combinatorics}
\LANGen{
\Question{
There are \(24\) 
girls and \(8\) 
boys in the class. How many ways are there to designate a president and vice-president
of the class if it is required that one of the position will be held by a boy and the
other one by a girl?}
{\(24\cdot 8\cdot 2=384\)}{\(24\cdot 8=192\)}{\(\frac{32!}
{2!\; 30!}=496\)}{\(\frac{32!}
{24!\; 8!}=10\:518\:300\)}{}{}}
\LANGpl{
\Question{
Klasa liczy \(24\) 
dziewczyny  i \(8\) 
chłopaków. Ile jest sposobów wybrania przewodniczącego klasy i wiceprzewodniczącego jeżeli jedną z tych funkcji musi mieć chłopak, a drugą dziewczyna?}
{\(24\cdot 8\cdot 2=384\)}{\(24\cdot 8=192\)}{\(\frac{32!}
{2!\; 30!}=496\)}{\(\frac{32!}
{24!\; 8!}=10\:518\:300\)}{}{}}
\LANGcs{
\Question{
Ve třídě je celkem \(24\) 
dívek a \(8\) 
chlapců. Určete, kolika způsoby můžeme vybrat předsedu a místopředsedu
třídy, jestliže jednu funkci bude zastávat dívka a druhou chlapec.}
{\(24\cdot 8\cdot 2=384\)}{\(24\cdot 8=192\)}{\(\frac{32!}
{2!\; 30!}=496\)}{\(\frac{32!}
{24!\; 8!}=10\:518\:300\)}{}{}}
\LANGsk{
\Question{
V triede je \(24\) 
dievčat a \(8\) 
chlapcov. Koľkými spôsobmi môžeme vybrať predsedu a podpredsedu triedy, ak jednu funkciu bude vykonávať dievča a druhú chlapec?}
{\(24\cdot 8\cdot 2=384\)}{\(24\cdot 8=192\)}{\(\frac{32!}
{2!\; 30!}=496\)}{\(\frac{32!}
{24!\; 8!}=10\:518\:300\)}{}{}}
\LANGes{
\Question{
Hay \(24\) 
chicas y \(8\) 
chicos en la clase. ¿De cuántas maneras se puede elegir un presidente y un vicepresidente
de la clase si se requiere que uno de los puestos está ocopado por un chico y el
otro por una chica?}
{\(24\cdot 8\cdot 2=384\)}{\(24\cdot 8=192\)}{\(\frac{32!}
{2!\; 30!}=496\)}{\(\frac{32!}
{24!\; 8!}=10\:518\:300\)}{}{}}
\End

\Data\ID{9000148910}
\Updated{2021-04-25 07:04:06}
\Level{A}
\SubArea{Combinatorics}
\LANGen{
\Question{
The player throws a dice in a desktop game. If he gets the number
\(6\), he
throws again and his score is the sum of both numbers. Find how many possibilities
are there for the player's score.}
{\(5 + 6=11\)}{\(2\cdot 6=12\)}{\(1 + 6=7\)}{\(6\cdot 6=36\)}{}{}}
\LANGpl{
\Question{
Gracz rzuca kostką do gry. Jeśli wyrzuci 
\(6\), rzuca ponownie, jego wynik to suma obu wyrzuconych liczb. Ile istnieje możliwych wyników jakie może uzyskać gracz?}
{\(5 + 6=11\)}{\(2\cdot 6=12\)}{\(1 + 6=7\)}{\(6\cdot 6=36\)}{}{}}
\LANGcs{
\Question{
Při hře „Člověče nezlob se” hází hráč šestistěnnou kostkou. Pokud
hodí šestku, hází ještě jednou. Pokud hodí šestku i podruhé,
potřetí už nehází. Kolika způsoby může hod/hody dopadnout?}
{\(5 + 6=11\)}{\(2\cdot 6=12\)}{\(1 + 6=7\)}{\(6\cdot 6=36\)}{}{}}
\LANGsk{
\Question{
Pri stolovej hre hráč hádže kockou. Ak hodí 
\(6\), hádže ešte raz. 
 Ak hodí šestku aj druhýkrát, tretíkrát už nehádže. Koľkými spôsobmi
môže hod dopadnúť?}
{\(5 + 6=11\)}{\(2\cdot 6=12\)}{\(1 + 6=7\)}{\(6\cdot 6=36\)}{}{}}
\LANGes{
\Question{
Un jugador lanza un dado en un juego de mesa. Si obtiene el número
\(6\),  lanza de nuevo y su puntuación es la suma de ambos números. Halla cuantas posibilidades
de puntuación existen para el jugador.}
{\(5 + 6=11\)}{\(2\cdot 6=12\)}{\(1 + 6=7\)}{\(6\cdot 6=36\)}{}{}}
\End

\Data\ID{9000153801}
\Updated{2021-04-25 08:04:03}
\Level{A}
\SubArea{Combinatorics}
\LANGen{
\Question{
Find the number of mutually different triangles such that each side of each triangle
is either \(2\),
\(3\),
\(4\) or
\(5\).}
{\(17\)}{\(14\)}{\(20\)}{\(16\)}{}{}}
\LANGpl{
\Question{
Wskaż liczbę różnych trójkątów tak, aby każdy bok trójkąta wynosił
 \(2\),
\(3\),
\(4\) lub
\(5\).}
{\(17\)}{\(14\)}{\(20\)}{\(16\)}{}{}}
\LANGcs{
\Question{
Určete počet všech trojúhelníků, z nichž žádné
dva nejsou shodné a jejichž každá strana má jednu
z velikostí daných čísly \(2\),
\(3\),
\(4\),
\(5\).}
{\(17\)}{\(14\)}{\(20\)}{\(16\)}{}{}}
\LANGsk{
\Question{
Určte počet všetkých trojuholníkov, z ktorých žiadne dva nie sú zhodné a ich každá strana má jednu z veľkostí daných číslami
 \(2\),
\(3\),
\(4\) alebo
\(5\).}
{\(17\)}{\(14\)}{\(20\)}{\(16\)}{}{}}
\LANGes{
\Question{
Determina el número de triángulos diferentes de modo que cada lado de cada triángulo
sea \(2\),
\(3\),
\(4\) o
\(5\).}
{\(17\)}{\(14\)}{\(20\)}{\(16\)}{}{}}
\End

\Data\ID{9000153802}
\Updated{2021-04-25 08:04:34}
\Level{A}
\SubArea{Combinatorics}
\LANGen{
\Question{
Find the number of mutually different isosceles triangles (at least
two sides are equal) such that each side of each triangle is either
\(2\),
\(3\),
\(4\) or
\(5\).}
{\(14\)}{\(16\)}{\(17\)}{\(24\)}{}{}}
\LANGpl{
\Question{
Wskaż liczbę różnych trójkątów równoramiennych tak, aby każdy bok trójkąta był równy
\(2\),
\(3\),
\(4\) lub
\(5\).}
{\(14\)}{\(16\)}{\(17\)}{\(24\)}{}{}}
\LANGcs{
\Question{
Určete počet všech rovnoramenných trojúhelníků, z nichž žádné dva
nejsou shodné a jejichž každá strana má jednu z velikostí daných čísly
\(2\),
\(3\),
\(4\),
\(5\).}
{\(14\)}{\(16\)}{\(17\)}{\(24\)}{}{}}
\LANGsk{
\Question{
Určte počet všetkých rovnoramenných trojuholníkov, z ktorých žiadne dva nie sú zhodné a každá ich strana má jednu z veľkostí daných číslami
\(2\),
\(3\),
\(4\) alebo
\(5\).}
{\(14\)}{\(16\)}{\(17\)}{\(24\)}{}{}}
\LANGes{
\Question{
Determina el número de triángulos isósceles diferentes (al menos
dos lados son iguales) de modo que cada lado de cada triángulo sea
\(2\),
\(3\),
\(4\) o
\(5\).}
{\(14\)}{\(16\)}{\(17\)}{\(24\)}{}{}}
\End

\Data\ID{9000153803}
\Updated{2021-04-25 08:04:13}
\Level{A}
\SubArea{Combinatorics}
\LANGen{
\Question{
Find the number of mutually different triangles such that all
three sides of each triangle are mutually different and each side is
\(2\),
\(3\),
\(4\) or
\(5\).}
{\(3\)}{\(4\)}{\(14\)}{\(16\)}{}{}}
\LANGpl{
\Question{
Wskaż liczbę różnych trójkątów tak, aby boki trójkąta były różne 
oraz każdy bok był równy
\(2\),
\(3\),
\(4\) lub
\(5\).}
{\(3\)}{\(4\)}{\(14\)}{\(16\)}{}{}}
\LANGcs{
\Question{
Určete počet všech nerovnoramenných trojúhelníků, z nichž žádné dva
nejsou shodné a jejichž každá strana má jednu z velikostí daných čísly
\(2\),
\(3\),
\(4\),
\(5\).}
{\(3\)}{\(4\)}{\(14\)}{\(16\)}{}{}}
\LANGsk{
\Question{
Určte počet všetkých nerovnoramenných trojuholníkov, z ktorých žiadne dva nie sú zhodné a každá ich strana má jednu z veľkostí daných číslami
 \(2\),
\(3\),
\(4\) alebo
\(5\).}
{\(3\)}{\(4\)}{\(14\)}{\(16\)}{}{}}
\LANGes{
\Question{
Determina el número de triángulos diferentes que se pueden formar de manera que todos
los lados de cada triángulo sean diferentes y cada lado sea
\(2\),
\(3\),
\(4\) o
\(5\).}
{\(3\)}{\(4\)}{\(14\)}{\(16\)}{}{}}
\End

\Data\ID{9000153804}
\Updated{2021-04-25 07:04:18}
\Level{A}
\SubArea{Combinatorics}
\LANGen{
\Question{
Find the number of \(3\)-combinations
from the set \(\{2,3,4,5\}\).}
{\(4\)}{\(6\)}{\(24\)}{\(20\)}{}{}}
\LANGpl{
\Question{
Wskaż liczbę kombinacji  \(3\)-elementówych
ze zbioru \(\{2,3,4,5\}\).}
{\(4\)}{\(6\)}{\(24\)}{\(20\)}{}{}}
\LANGcs{
\Question{
Určete počet tříčlenných kombinací prvků množiny
\(\{2,3,4,5\}\).}
{\(4\)}{\(6\)}{\(24\)}{\(20\)}{}{}}
\LANGsk{
\Question{
Určte počet  \(3\)-členných kombinácií prvkov množiny 
 \(\{2,3,4,5\}\).}
{\(4\)}{\(6\)}{\(24\)}{\(20\)}{}{}}
\LANGes{
\Question{
Halla el número de combinaciones de \(3\) elementos
del conjunto \(\{2,3,4,5\}\).}
{\(4\)}{\(6\)}{\(24\)}{\(20\)}{}{}}
\End

\Data\ID{9000153805}
\Updated{2021-04-25 07:04:04}
\Level{A}
\SubArea{Combinatorics}
\LANGen{
\Question{
Find the number of \(3\)-combinations
with repetition from the set \(\{2,3,4,5\}\).}
{\(20\)}{\(4\)}{\(24\)}{\(12\)}{}{}}
\LANGpl{
\Question{
Wskaż liczbę kombinacji  \(3\)-elementówych
z powtórzeniami ze zbioru \(\{2,3,4,5\}\).}
{\(20\)}{\(4\)}{\(24\)}{\(12\)}{}{}}
\LANGcs{
\Question{
Určete počet tříčlenných kombinací s opakováním prvků množiny
\(\{2,3,4,5\}\).}
{\(20\)}{\(4\)}{\(24\)}{\(12\)}{}{}}
\LANGsk{
\Question{
Určte počet \(3\)-členných kombinácií
s opakovaním prvkov množiny \(\{2,3,4,5\}\).}
{\(20\)}{\(4\)}{\(24\)}{\(12\)}{}{}}
\LANGes{
\Question{
Halla el número de combinaciones de \(3\) elementos con repetición del conjunto \(\{2,3,4,5\}\).}
{\(20\)}{\(4\)}{\(24\)}{\(12\)}{}{}}
\End

\Data\ID{9000153806}
\Updated{2020-08-19 01:08:42}
\Level{A}
\SubArea{Combinatorics}
\LANGen{
\Question{
Find the number of the positive integers with three digits which can be written using
the digits \(2\),
\(3\),
\(4\) and
\(5\). The
digits may be reused.}
{\(64\)}{\(24\)}{\(256\)}{\(81\)}{}{}}
\LANGpl{
\Question{
Wskaż liczbę trzycyfrowych  dodatnich liczb całkowitych, które mogą być zapisane za pomocą 
cyfr \(2\),
\(3\),
\(4\) i
\(5\). Cyfry mogą zostać użyte ponownie.}
{\(64\)}{\(24\)}{\(256\)}{\(81\)}{}{}}
\LANGcs{
\Question{
Určete počet trojciferných přirozených čísel, jež lze sestavit pouze z cifer
\(2\),
\(3\),
\(4\),
\(5\).}
{\(64\)}{\(24\)}{\(256\)}{\(81\)}{}{}}
\LANGsk{
\Question{
Určte počet trojciferných prirodzených čísel, ktoré možno zostaviť len z cifier 
 \(2\),
\(3\),
\(4\) a
\(5\).}
{\(64\)}{\(24\)}{\(256\)}{\(81\)}{}{}}
\LANGes{
\Question{
Determina el número de enteros positivos con tres dígitos que se pueden escribir usando
los dígitos \(2\),
\(3\),
\(4\) y
\(5\). Los dígitos pueden reutilizarse.}
{\(64\)}{\(24\)}{\(256\)}{\(81\)}{}{}}
\End

\Data\ID{9000153807}
\Updated{2020-08-19 01:08:05}
\Level{A}
\SubArea{Combinatorics}
\LANGen{
\Question{
Find the number of the positive integers with three digits which can be written using
the digits \(2\),
\(3\),
\(4\) and
\(5\). Each
digit can be used at most once.}
{\(24\)}{\(64\)}{\(256\)}{\(81\)}{}{}}
\LANGpl{
\Question{
Wskaż liczbę trzycyfrowych dodatnich liczb całkowitych, które mogą być zapisane za pomocą cyfr  \(2\),
\(3\),
\(4\) i
\(5\). Każda cyfra może być użyta tylko raz.}
{\(24\)}{\(64\)}{\(256\)}{\(81\)}{}{}}
\LANGcs{
\Question{
Určete počet trojciferných přirozených čísel s různými ciframi, jež lze sestavit
pouze z číslic \(2\),
\(3\),
\(4\),
\(5\).}
{\(24\)}{\(64\)}{\(256\)}{\(81\)}{}{}}
\LANGsk{
\Question{
Určte počet trojciferných prirodzených čísel s rôznymi ciframi, ktoré možno zostaviť len z číslic \(2\),
\(3\),
\(4\) a
\(5\). Každá
číslica môže byť použitá maximálne raz.}
{\(24\)}{\(64\)}{\(256\)}{\(81\)}{}{}}
\LANGes{
\Question{
Determina el número de enteros positivos con tres dígitos que se puede escribir usando
los dígitos  \(2\),
\(3\),
\(4\) y
\(5\). Cada dígito se puede utilizar como máximo una vez.}
{\(24\)}{\(64\)}{\(256\)}{\(81\)}{}{}}
\End

\Data\ID{9000153808}
\Updated{2020-08-19 01:08:23}
\Level{A}
\SubArea{Combinatorics}
\LANGen{
\Question{
Find the number of subsets of the set
\(\{2,3,4,5\}\).}
{\(16\)}{\(4\)}{\(6\)}{\(15\)}{}{}}
\LANGpl{
\Question{
Wskaż liczbę podzbiorów zbioru
\(\{2,3,4,5\}\).}
{\(16\)}{\(4\)}{\(6\)}{\(15\)}{}{}}
\LANGcs{
\Question{
Určete počet všech podmnožin množiny
\(\{2,3,4,5\}\).}
{\(16\)}{\(4\)}{\(6\)}{\(15\)}{}{}}
\LANGsk{
\Question{
Určte počet všetkých podmnožín množiny
\(\{2,3,4,5\}\).}
{\(16\)}{\(4\)}{\(6\)}{\(15\)}{}{}}
\LANGes{
\Question{
Determina el número de subconjuntos del conjunto
\(\{2,3,4,5\}\).}
{\(16\)}{\(4\)}{\(6\)}{\(15\)}{}{}}
\End

\Data\ID{9000153809}
\Updated{2021-04-25 07:04:00}
\Level{A}
\SubArea{Combinatorics}
\LANGen{
\Question{
Find the number of positive integers with three mutually different digits which can be written just
using the digits \(2\),
\(3\),
\(4\),
\(5\) and which can
be divided by \(4\).}
{\(6\)}{\(9\)}{\(10\)}{\(5\)}{}{}}
\LANGpl{
\Question{
Wskaż liczbę  dodatnich liczb całkowitych o trzech różnych cyfrach, które mogą być zapisane tylko za pomocą cyfr  \(2\),
\(3\),
\(4\),
\(5\) oraz muszą być podzielne przez \(4\).}
{\(6\)}{\(9\)}{\(10\)}{\(5\)}{}{}}
\LANGcs{
\Question{
Určete počet trojciferných přirozených čísel s různými ciframi, jež lze sestavit
pouze z číslic \(2\),
\(3\),
\(4\),
\(5\) 
a která jsou dělitelná čtyřmi.}
{\(6\)}{\(9\)}{\(10\)}{\(5\)}{}{}}
\LANGsk{
\Question{
Určte počet trojciferných prirodzených čísel s rôznymi ciframi, ktoré možno zostaviť len z číslic \(2\),
\(3\),
\(4\),
\(5\) a ktoré sú deliteľné \(4\).}
{\(6\)}{\(9\)}{\(10\)}{\(5\)}{}{}}
\LANGes{
\Question{
Determina el número de enteros positivos con tres dígitos diferentes que se pueden escribir simplemente
usando los dígitos \(2\),
\(3\), 
\(4\),
\(5\) y son divisibles
 por \(4\).}
{\(6\)}{\(9\)}{\(10\)}{\(5\)}{}{}}
\End

\Data\ID{9000153810}
\Updated{2021-04-25 07:04:10}
\Level{A}
\SubArea{Combinatorics}
\LANGen{
\Question{
Find the number of positive integers with three mutually different digits which can be written just
using the digits \(2\),
\(3\),
\(4\),
\(5\) and which can
be divided by \(3\).}
{\(12\)}{\(6\)}{\(9\)}{\(10\)}{}{}}
\LANGpl{
\Question{
Wskaż liczbę dodatnich liczb całkowitych o różnych cyfrach, które mogą być zapisane tylko za pomocą  \(2\),
\(3\),
\(4\),
\(5\) oraz muszą być podzielne przez  \(3\).}
{\(12\)}{\(6\)}{\(9\)}{\(10\)}{}{}}
\LANGcs{
\Question{
Určete počet trojciferných přirozených čísel s různými ciframi, jež lze sestavit
pouze z číslic \(2\),
\(3\),
\(4\),
\(5\) 
a která jsou dělitelná třemi.}
{\(12\)}{\(6\)}{\(9\)}{\(10\)}{}{}}
\LANGsk{
\Question{
Určte počet trojciferných prirodzených čísel s rôznymi ciframi, ktoré možno zostaviť len z číslic 
 \(2\),
\(3\),
\(4\),
\(5\) a ktoré sú deliteľné \(3\).}
{\(12\)}{\(6\)}{\(9\)}{\(10\)}{}{}}
\LANGes{
\Question{
Determina el número de enteros positivos con tres dígitos diferentes que se pueden escribir simplemente usando los dígitos 
 \(2\),
\(3\),
\(4\),
\(5\) y son divisibles por \(3\).}
{\(12\)}{\(6\)}{\(9\)}{\(10\)}{}{}}
\End

\Data\ID{1003024603}
\Updated{2021-04-25 06:04:13}
\Level{B}
\SubArea{Combinatorics}
\LANGen{
\Question{
Simplify the following expression without using a calculator: 
\( \frac{85!+4\cdot84!-84\cdot83!}{2\cdot84!}\)}
{\( 44 \)}{\( 88 \)}{\( 66 \)}{\( 22 \)}{}{}}
\LANGpl{
\Question{
Uprość bez korzystania z kalkulatora: 
\( \frac{85!+4\cdot84!-84\cdot83!}{2\cdot84!}\)}
{\( 44 \)}{\( 88 \)}{\( 66 \)}{\( 22 \)}{}{}}
\LANGcs{
\Question{
Zjednodušte bez použití kalkulačky: 
\( \frac{85!+4\cdot84!-84\cdot83!}{2\cdot84!}\)}
{\( 44 \)}{\( 88 \)}{\( 66 \)}{\( 22 \)}{}{}}
\LANGsk{
\Question{
Zjednodušte bez použitia kalkulačky: 
\( \frac{85!+4\cdot84!-84\cdot83!}{2\cdot84!}\)}
{\( 44 \)}{\( 88 \)}{\( 66 \)}{\( 22 \)}{}{}}
\LANGes{
\Question{
Simplifica la siguiente expresión sin usar calculadora:
\( \frac{85!+4\cdot84!-84\cdot83!}{2\cdot84!}\)}
{\( 44 \)}{\( 88 \)}{\( 66 \)}{\( 22 \)}{}{}}
\End

\Data\ID{1003024604}
\Updated{2021-04-25 06:04:33}
\Level{B}
\SubArea{Combinatorics}
\LANGen{
\Question{
Without using a calculator, identify which of the following values is the greatest one.}
{the number of ordered arrangements of \( 5 \) objects taking \( 3 \) at a time}{the number of unordered selections  of \( 3 \) objects  out of \( 5 \) different kinds  with repetition allowed}{the number of ordered arrangements of \( 5 \) objects in which \( 3 \) are identical}{the number of unordered  selections of \( 3 \) objects out of \( 5 \) objects}{}{}}
\LANGpl{
\Question{
Wskaż, która z wartości jest największa. Nie używaj kalkulatora.}
{wariacja trzech elementów z pięciu}{kombinacja trzech elementów z pięciu z permutacją z powtórzeniami}{permutacja z powtórzeniami, w których pomiędzy pięcioma elementami jeden powtarza się 3 razy}{kombinacja   trzech elementów z pięciu}{}{}}
\LANGcs{
\Question{
Bez použití kalkulačky určete, které z čísel je největší.}
{počet variací třetí třídy z pěti prvků}{počet kombinací třetí třídy z pěti prvků s opakováním}{počet permutací s opakováním, kde se mezi pěti prvky jeden opakuje 3 krát}{počet kombinací třetí třídy z pěti prvků}{}{}}
\LANGsk{
\Question{
Bez použitia kalkulačky určite, ktoré z čísel je najväčšie.}
{počet variácií tretej triedy z piatich prvkov}{počet kombinácií tretej triedy z piatich prvkov s opakovaním}{počet permutácií s opakovaním, kde sa medzi piatimi prvkami jeden opakuje 3 krát}{počet kombinácií tretej triedy z piatich prvkov}{}{}}
\LANGes{
\Question{
Sin usar calculadora, identifica cuál de los siguientes valores es el mayor.}
{El número de variaciones de tercer orden de 5 elementos.}{El número de combinaciones de tercer orden de 5 elementos con repetición.}{El número de permutaciones con repetición, donde uno de los cinco elementos se repite 3 veces.}{El número de combinaciones de tercer orden de 5 elementos.}{}{}}
\End

\Data\ID{1003024605}
\Updated{2021-04-25 06:04:00}
\Level{B}
\SubArea{Combinatorics}
\LANGen{
\Question{
Simplify the following expression without using a calculator:
\( \frac{91}{13!}+\frac{6}{12!}-\frac1{11!} \)}
{\( \frac1{12!} \)}{\( \frac{96}{13!} \)}{\( \frac{85}{13!} \)}{\( \frac1{13!} \)}{}{}}
\LANGpl{
\Question{
Uprość bez korzystania z kalkulatora:
\( \frac{91}{13!}+\frac{6}{12!}-\frac1{11!} \)}
{\( \frac1{12!} \)}{\( \frac{96}{13!} \)}{\( \frac{85}{13!} \)}{\( \frac1{13!} \)}{}{}}
\LANGcs{
\Question{
Zjednodušte bez použití kalkulačky:
\( \frac{91}{13!}+\frac{6}{12!}-\frac1{11!} \)}
{\( \frac1{12!} \)}{\( \frac{96}{13!} \)}{\( \frac{85}{13!} \)}{\( \frac1{13!} \)}{}{}}
\LANGsk{
\Question{
Zjednodušte bez použitia kalkulačky:
\( \frac{91}{13!}+\frac{6}{12!}-\frac1{11!} \)}
{\( \frac1{12!} \)}{\( \frac{96}{13!} \)}{\( \frac{85}{13!} \)}{\( \frac1{13!} \)}{}{}}
\LANGes{
\Question{
Simplifica la siguiente expresión sin usar calculadora:
\( \frac{91}{13!}+\frac{6}{12!}-\frac1{11!} \)}
{\( \frac1{12!} \)}{\( \frac{96}{13!} \)}{\( \frac{85}{13!} \)}{\( \frac1{13!} \)}{}{}}
\End

\Data\ID{9000136901}
\Updated{2021-07-02 09:07:48}
\Level{B}
\SubArea{Combinatorics}
\LANGen{
\Question{
Simplify \(\left({15\above 0.0pt 
 8} \right) +\left ({15\above 0.0pt 
 9} \right)\).}
{\(\left({16\above 0.0pt 
 9} \right)\)}{\(\left({15\above 0.0pt 
10}\right)\)}{\(\left({15\above 0.0pt 
 7} \right)\)}{\(\left({16\above 0.0pt 
 8} \right)\)}{\(\left({30\above 0.0pt 
17}\right)\)}{}}
\LANGpl{
\Question{
Uprość  \(\left({15\above 0.0pt 
 8} \right) +\left ({15\above 0.0pt 
 9} \right)\).}
{\(\left({16\above 0.0pt 
 9} \right)\)}{\(\left({15\above 0.0pt 
10}\right)\)}{\(\left({15\above 0.0pt 
 7} \right)\)}{\(\left({16\above 0.0pt 
 8} \right)\)}{\(\left({30\above 0.0pt 
17}\right)\)}{}}
\LANGcs{
\Question{
Součet \(\left({15\above 0.0pt 
 8} \right) +\left ({15\above 0.0pt 
 9} \right)\) 
je roven:}
{\(\left({16\above 0.0pt 
 9} \right)\)}{\(\left({15\above 0.0pt 
10}\right)\)}{\(\left({15\above 0.0pt 
 7} \right)\)}{\(\left({16\above 0.0pt 
 8} \right)\)}{\(\left({30\above 0.0pt 
17}\right)\)}{}}
\LANGsk{
\Question{
Zjednodušte \(\left({15\above 0.0pt 
 8} \right) +\left ({15\above 0.0pt 
 9} \right)\).}
{\(\left({16\above 0.0pt 
 9} \right)\)}{\(\left({15\above 0.0pt 
10}\right)\)}{\(\left({15\above 0.0pt 
 7} \right)\)}{\(\left({16\above 0.0pt 
 8} \right)\)}{\(\left({30\above 0.0pt 
17}\right)\)}{}}
\LANGes{
\Question{
Simplifica \(\left({15\above 0.0pt 
 8} \right) +\left ({15\above 0.0pt 
 9} \right)\).}
{\(\left({16\above 0.0pt 
 9} \right)\)}{\(\left({15\above 0.0pt 
10}\right)\)}{\(\left({15\above 0.0pt 
 7} \right)\)}{\(\left({16\above 0.0pt 
 8} \right)\)}{\(\left({30\above 0.0pt 
17}\right)\)}{}}
\End

\Data\ID{9000136902}
\Updated{2020-07-19 07:07:37}
\Level{B}
\SubArea{Combinatorics}
\LANGen{
\Question{
Simplify \(\left({12\above 0.0pt 
10}\right) -\left ({12\above 0.0pt 
 2} \right)\).}
{\(0\)}{\(66\)}{\(\left({12\above 0.0pt 
 8} \right)\)}{\(1\)}{\(\left({12\above 0.0pt 
 0} \right)\)}{}}
\LANGpl{
\Question{
Uprość  \(\left({12\above 0.0pt 
10}\right) -\left ({12\above 0.0pt 
 2} \right)\).}
{\(0\)}{\(66\)}{\(\left({12\above 0.0pt 
 8} \right)\)}{\(1\)}{\(\left({12\above 0.0pt 
 0} \right)\)}{}}
\LANGcs{
\Question{
Rozdíl \(\left({12\above 0.0pt 
10}\right) -\left ({12\above 0.0pt 
 2} \right)\) 
je roven:}
{\(0\)}{\(66\)}{\(\left({12\above 0.0pt 
 8} \right)\)}{\(1\)}{\(\left({12\above 0.0pt 
 0} \right)\)}{}}
\LANGsk{
\Question{
Zjednodušte \(\left({12\above 0.0pt 
10}\right) -\left ({12\above 0.0pt 
 2} \right)\).}
{\(0\)}{\(66\)}{\(\left({12\above 0.0pt 
 8} \right)\)}{\(1\)}{\(\left({12\above 0.0pt 
 0} \right)\)}{}}
\LANGes{
\Question{
Simplifica \(\left({12\above 0.0pt 
10}\right) -\left ({12\above 0.0pt 
 2} \right)\).}
{\(0\)}{\(66\)}{\(\left({12\above 0.0pt 
 8} \right)\)}{\(1\)}{\(\left({12\above 0.0pt 
 0} \right)\)}{}}
\End

\Data\ID{9000136903}
\Updated{2020-07-19 07:07:17}
\Level{B}
\SubArea{Combinatorics}
\LANGen{
\Question{
Simplify \(\left({4\above 0.0pt 
0}\right) +\left ({4\above 0.0pt 
1}\right) +\left ({4\above 0.0pt 
2}\right) +\left ({4\above 0.0pt 
3}\right) +\left ({4\above 0.0pt 
4}\right)\).}
{\(4^{2}\)}{\(14\)}{\(\left({5\above 0.0pt 
4}\right)\)}{\(32\)}{\(\left({8\above 0.0pt 
4}\right)\)}{}}
\LANGpl{
\Question{
Uprość  \(\left({4\above 0.0pt 
0}\right) +\left ({4\above 0.0pt 
1}\right) +\left ({4\above 0.0pt 
2}\right) +\left ({4\above 0.0pt 
3}\right) +\left ({4\above 0.0pt 
4}\right)\).}
{\(4^{2}\)}{\(14\)}{\(\left({5\above 0.0pt 
4}\right)\)}{\(32\)}{\(\left({8\above 0.0pt 
4}\right)\)}{}}
\LANGcs{
\Question{
Součet \(\left({4\above 0.0pt 
0}\right) +\left ({4\above 0.0pt 
1}\right) +\left ({4\above 0.0pt 
2}\right) +\left ({4\above 0.0pt 
3}\right) +\left ({4\above 0.0pt 
4}\right)\) 
je roven:}
{\(4^{2}\)}{\(14\)}{\(\left({5\above 0.0pt 
4}\right)\)}{\(32\)}{\(\left({8\above 0.0pt 
4}\right)\)}{}}
\LANGsk{
\Question{
Zjednodušte \(\left({4\above 0.0pt 
0}\right) +\left ({4\above 0.0pt 
1}\right) +\left ({4\above 0.0pt 
2}\right) +\left ({4\above 0.0pt 
3}\right) +\left ({4\above 0.0pt 
4}\right)\).}
{\(4^{2}\)}{\(14\)}{\(\left({5\above 0.0pt 
4}\right)\)}{\(32\)}{\(\left({8\above 0.0pt 
4}\right)\)}{}}
\LANGes{
\Question{
Simplifica \(\left({4\above 0.0pt 
0}\right) +\left ({4\above 0.0pt 
1}\right) +\left ({4\above 0.0pt 
2}\right) +\left ({4\above 0.0pt 
3}\right) +\left ({4\above 0.0pt 
4}\right)\).}
{\(4^{2}\)}{\(14\)}{\(\left({5\above 0.0pt 
4}\right)\)}{\(32\)}{\(\left({8\above 0.0pt 
4}\right)\)}{}}
\End

\Data\ID{9000136904}
\Updated{2020-07-19 07:07:46}
\Level{B}
\SubArea{Combinatorics}
\LANGen{
\Question{
Simplify \(\left({n+1\above 0.0pt 
 n} \right) +\left ({n+1\above 0.0pt 
 1} \right)\) 
for \(n\in \mathbb{N}\).}
{\(2(n + 1)\)}{\(n + 2\)}{\(\left({n+2\above 0.0pt 
 n} \right)\)}{\(2\)}{\(\left (n + 1\right )^{2}\)}{}}
\LANGpl{
\Question{
Uprość  \(\left({n+1\above 0.0pt 
 n} \right) +\left ({n+1\above 0.0pt 
 1} \right)\) 
dla  \(n\in \mathbb{N}\).}
{\(2(n + 1)\)}{\(n + 2\)}{\(\left({n+2\above 0.0pt 
 n} \right)\)}{\(2\)}{\(\left (n + 1\right )^{2}\)}{}}
\LANGcs{
\Question{
Součet \(\left({n+1\above 0.0pt 
 n} \right) +\left ({n+1\above 0.0pt 
 1} \right)\) je
pro libovolné \(n\in \mathbb{N}\) 
roven:}
{\(2(n + 1)\)}{\(n + 2\)}{\(\left({n+2\above 0.0pt 
 n} \right)\)}{\(2\)}{\(\left (n + 1\right )^{2}\)}{}}
\LANGsk{
\Question{
Zjednodušte \(\left({n+1\above 0.0pt 
 n} \right) +\left ({n+1\above 0.0pt 
 1} \right)\) 
pre \(n\in \mathbb{N}\).}
{\(2(n + 1)\)}{\(n + 2\)}{\(\left({n+2\above 0.0pt 
 n} \right)\)}{\(2\)}{\(\left (n + 1\right )^{2}\)}{}}
\LANGes{
\Question{
Simplifica \(\left({n+1\above 0.0pt 
 n} \right) +\left ({n+1\above 0.0pt 
 1} \right)\) 
para \(n\in \mathbb{N}\).}
{\(2(n + 1)\)}{\(n + 2\)}{\(\left({n+2\above 0.0pt 
 n} \right)\)}{\(2\)}{\(\left (n + 1\right )^{2}\)}{}}
\End

\Data\ID{9000136905}
\Updated{2020-07-19 07:07:03}
\Level{B}
\SubArea{Combinatorics}
\LANGen{
\Question{
Simplify \(\left({n\above 0.0pt 
2} \right) -\left ({ n\above 0.0pt 
n-2}\right)\) 
for \(n\in \mathbb{N}\),
\(n\geq 2\).}
{\(0\)}{\(\left (n + 2\right )\left (n + 1\right )\)}{\(\left({n+2\above 0.0pt 
 n} \right)\)}{\(n^{2} - 1\)}{\(\left({n\above 0.0pt 
n}\right)\)}{}}
\LANGpl{
\Question{
Uprość \(\left({n\above 0.0pt 
2} \right) -\left ({ n\above 0.0pt 
n-2}\right)\) 
dla  \(n\in \mathbb{N}\),
\(n\geq 2\).}
{\(0\)}{\(\left (n + 2\right )\left (n + 1\right )\)}{\(\left({n+2\above 0.0pt 
 n} \right)\)}{\(n^{2} - 1\)}{\(\left({n\above 0.0pt 
n}\right)\)}{}}
\LANGcs{
\Question{
Rozdíl \(\left({n\above 0.0pt 
2} \right) -\left ({ n\above 0.0pt 
n-2}\right)\) je pro
libovolné \(n\in \mathbb{N}\),
\(n\geq 2\),
roven:}
{\(0\)}{\(\left (n + 2\right )\left (n + 1\right )\)}{\(\left({n+2\above 0.0pt 
 n} \right)\)}{\(n^{2} - 1\)}{\(\left({n\above 0.0pt 
n}\right)\)}{}}
\LANGsk{
\Question{
Zjednodušte \(\left({n\above 0.0pt 
2} \right) -\left ({ n\above 0.0pt 
n-2}\right)\) 
pre \(n\in \mathbb{N}\),
\(n\geq 2\).}
{\(0\)}{\(\left (n + 2\right )\left (n + 1\right )\)}{\(\left({n+2\above 0.0pt 
 n} \right)\)}{\(n^{2} - 1\)}{\(\left({n\above 0.0pt 
n}\right)\)}{}}
\LANGes{
\Question{
Simplifica \(\left({n\above 0.0pt 
2} \right) -\left ({ n\above 0.0pt 
n-2}\right)\) 
para \(n\in \mathbb{N}\),
\(n\geq 2\).}
{\(0\)}{\(\left (n + 2\right )\left (n + 1\right )\)}{\(\left({n+2\above 0.0pt 
 n} \right)\)}{\(n^{2} - 1\)}{\(\left({n\above 0.0pt 
n}\right)\)}{}}
\End

\Data\ID{9000136906}
\Updated{2021-04-25 06:04:56}
\Level{B}
\SubArea{Combinatorics}
\LANGen{
\Question{
Solve the equation involving binomial coefficient. 
\[ 
 \left({x\above 0.0pt 
2}\right) = 15
\]}
{\(6\)}{\(5\)}{\(4\)}{\(3\)}{The equation has no solution.}{}}
\LANGpl{
\Question{
Rozwiąż równanie z symbolem Newtona.
\[ 
 \left({x\above 0.0pt 
2}\right) = 15
\]}
{\(6\)}{\(5\)}{\(4\)}{\(3\)}{Równanie nie ma rozwiązania.}{}}
\LANGcs{
\Question{
Vyberte celé číslo, které je řešením rovnice
\(\left({x\above 0.0pt 
2} \right) = 15\).
Pokud takové celé číslo neexistuje, tak zaškrtněte, že rovnice nemá
řešení.}
{\(6\)}{\(5\)}{\(4\)}{\(3\)}{Rovnice nemá řešení.}{}}
\LANGsk{
\Question{
Vyberte celé číslo, ktoré je riešením danej rovnice. Ak také číslo neexistuje, tak označte, že rovnica nemá riešenie.
\[ 
 \left({x\above 0.0pt 
2}\right) = 15
\]}
{\(6\)}{\(5\)}{\(4\)}{\(3\)}{Rovnica nemá riešenie.}{}}
\LANGes{
\Question{
Selecciona un número entero que sea solución de la ecuación
\[ 
 \left({x\above 0.0pt 
2}\right) = 15
\]}
{\(6\)}{\(5\)}{\(4\)}{\(3\)}{La ecuación no tiene ninguna solución.}{}}
\End

\Data\ID{9000136907}
\Updated{2021-04-25 06:04:29}
\Level{B}
\SubArea{Combinatorics}
\LANGen{
\Question{
Solve the equation involving binomial coefficients. 
\[ 
 \left({x + 1\above 0.0pt 
 x} \right) -\left ({x + 1\above 0.0pt 
x + 1}\right) = 21
\]}
{\(21\)}{\(20\)}{\(11\)}{\(10\)}{The equation has no solution.}{}}
\LANGpl{
\Question{
Rozwiąż równanie z symbolem Newtona.
\[ 
 \left({x + 1\above 0.0pt 
 x} \right) -\left ({x + 1\above 0.0pt 
x + 1}\right) = 21
\]}
{\(21\)}{\(20\)}{\(11\)}{\(10\)}{Równanie nie ma rozwiązania.}{}}
\LANGcs{
\Question{
Vyberte celé číslo, které je řešením rovnice
\(\left({x+1\above 0.0pt 
 x} \right) -\left ({x+1\above 0.0pt 
x+1}\right) = 21\).
Pokud takové celé číslo neexistuje, tak zaškrtněte, že rovnice nemá
řešení.}
{\(21\)}{\(20\)}{\(11\)}{\(10\)}{Rovnice nemá řešení.}{}}
\LANGsk{
\Question{
Vyberte celé číslo, ktoré je riešením danej rovnice. Ak také číslo neexistuje, tak označte, že rovnica nemá riešenie.
\[ 
 \left({x + 1\above 0.0pt 
 x} \right) -\left ({x + 1\above 0.0pt 
x + 1}\right) = 21
\]}
{\(21\)}{\(20\)}{\(11\)}{\(10\)}{Rovnica nemá riešenie.}{}}
\LANGes{
\Question{
Selecciona un número entero que sea solución de la ecuación
\[ 
 \left({x + 1\above 0.0pt 
 x} \right) -\left ({x + 1\above 0.0pt 
x + 1}\right) = 21
\]}
{\(21\)}{\(20\)}{\(11\)}{\(10\)}{La ecuación no tiene ninguna solución.}{}}
\End

\Data\ID{9000136908}
\Updated{2021-04-25 06:04:18}
\Level{B}
\SubArea{Combinatorics}
\LANGen{
\Question{
Solve the equation involving binomial coefficients. 
\[ 
 \left({x + 1\above 0.0pt 
 x} \right) -\left ({x + 1\above 0.0pt 
 1} \right) = 1
\]}
{The equation has no solution.}{\(- 1\)}{\(1\)}{\(2\)}{\(0\)}{}}
\LANGpl{
\Question{
Rozwiąż równanie z symbolem Newtona.
\[ 
 \left({x + 1\above 0.0pt 
 x} \right) -\left ({x + 1\above 0.0pt 
 1} \right) = 1
\]}
{Równanie nie ma rozwiązania.}{\(- 1\)}{\(1\)}{\(2\)}{\(0\)}{}}
\LANGcs{
\Question{
Vyberte celé číslo, které je řešením rovnice
\(\left({x+1\above 0.0pt 
 x} \right) -\left ({x+1\above 0.0pt 
 1} \right) = 1\).
Pokud takové celé číslo neexistuje, tak zaškrtněte, že rovnice nemá
řešení.}
{Rovnice nemá řešení.}{\(- 1\)}{\(1\)}{\(2\)}{\(0\)}{}}
\LANGsk{
\Question{
Vyberte celé číslo, ktoré je riešením danej rovnice. Ak také číslo neexistuje, tak označte, že rovnica nemá riešenie.
\[ 
 \left({x + 1\above 0.0pt 
 x} \right) -\left ({x + 1\above 0.0pt 
 1} \right) = 1
\]}
{Rovnica nemá riešenie.}{\(- 1\)}{\(1\)}{\(2\)}{\(0\)}{}}
\LANGes{
\Question{
Selecciona un número entero que sea solución de la ecuación
\[ 
 \left({x + 1\above 0.0pt 
 x} \right) -\left ({x + 1\above 0.0pt 
 1} \right) = 1
\]}
{La ecuación no tiene ninguna solución.}{\(- 1\)}{\(1\)}{\(2\)}{\(0\)}{}}
\End

\Data\ID{9000136909}
\Updated{2021-04-25 06:04:41}
\Level{B}
\SubArea{Combinatorics}
\LANGen{
\Question{
Solve the equation involving binomial coefficients. 
\[ 
 \left({x + 1\above 0.0pt 
 x} \right) +\left ({x + 2\above 0.0pt 
x + 1}\right) = 19
\]}
{\(8\)}{\(10\)}{\(12\)}{\(19\)}{The equation has no solution.}{}}
\LANGpl{
\Question{
Rozwiąż równanie z symbolem Newtona.
\[ 
 \left({x + 1\above 0.0pt 
 x} \right) +\left ({x + 2\above 0.0pt 
x + 1}\right) = 19
\]}
{\(8\)}{\(10\)}{\(12\)}{\(19\)}{Równanie nie ma rozwiązania.}{}}
\LANGcs{
\Question{
Vyberte celé číslo, které je řešením rovnice
\(\left({x+1\above 0.0pt 
 x} \right) +\left ({x+2\above 0.0pt 
x+1}\right) = 19\).
Pokud takové celé číslo neexistuje, tak zaškrtněte, že rovnice nemá
řešení.}
{\(8\)}{\(10\)}{\(12\)}{\(19\)}{Rovnice nemá řešení.}{}}
\LANGsk{
\Question{
Vyberte celé číslo, ktoré je riešením danej rovnice. Ak také číslo neexistuje, tak označte, že rovnica nemá riešenie. 
\[ 
 \left({x + 1\above 0.0pt 
 x} \right) +\left ({x + 2\above 0.0pt 
x + 1}\right) = 19
\]}
{\(8\)}{\(10\)}{\(12\)}{\(19\)}{Rovnica nemá riešenie.}{}}
\LANGes{
\Question{
Selecciona un número entero que sea solución de la ecuación
\[ 
 \left({x + 1\above 0.0pt 
 x} \right) +\left ({x + 2\above 0.0pt 
x + 1}\right) = 19
\]}
{\(8\)}{\(10\)}{\(12\)}{\(19\)}{La ecuación no tiene ninguna solución.}{}}
\End

\Data\ID{9000136910}
\Updated{2021-04-25 06:04:05}
\Level{B}
\SubArea{Combinatorics}
\LANGen{
\Question{
Solve the equation involving binomial coefficients. 
\[ 
 \left({x\above 0.0pt 
0}\right) +\left ({x\above 0.0pt 
1}\right) +\left ({x + 1\above 0.0pt 
 x} \right) = 25
\]}
{The equation has no solution.}{\(9\)}{\(1\)}{\(5\)}{\(7\)}{}}
\LANGpl{
\Question{
Rozwiąż równanie z symbolem Newtona.
\[ 
 \left({x\above 0.0pt 
0}\right) +\left ({x\above 0.0pt 
1}\right) +\left ({x + 1\above 0.0pt 
 x} \right) = 25
\]}
{Równanie nie ma rozwiązania.}{\(9\)}{\(1\)}{\(5\)}{\(7\)}{}}
\LANGcs{
\Question{
Vyberte celé číslo, které je řešením rovnice
\(\left({x\above 0.0pt 
0} \right) +\left ({x\above 0.0pt 
1} \right) +\left ({x+1\above 0.0pt 
 x} \right) = 25\).
Pokud takové celé číslo neexistuje, tak zaškrtněte, že rovnice nemá
řešení.}
{Rovnice nemá řešení.}{\(9\)}{\(1\)}{\(5\)}{\(7\)}{}}
\LANGsk{
\Question{
Vyberte celé číslo, ktoré je riešením danej rovnice. Ak také číslo neexistuje, tak označte, že rovnica nemá riešenie. 
\[ 
 \left({x\above 0.0pt 
0}\right) +\left ({x\above 0.0pt 
1}\right) +\left ({x + 1\above 0.0pt 
 x} \right) = 25
\]}
{Rovnica nemá riešenie.}{\(9\)}{\(1\)}{\(5\)}{\(7\)}{}}
\LANGes{
\Question{
Selecciona un número entero que sea solución de la ecuación
\[ 
 \left({x\above 0.0pt 
0}\right) +\left ({x\above 0.0pt 
1}\right) +\left ({x + 1\above 0.0pt 
 x} \right) = 25
\]}
{La ecuación no tiene ninguna solución.}{\(9\)}{\(1\)}{\(5\)}{\(7\)}{}}
\End

\Data\ID{9000140501}
\Updated{2020-08-19 11:08:59}
\Level{B}
\SubArea{Combinatorics}
\LANGen{
\Question{
Simplify for \(n\in \mathbb{N}\).
\[ 
 \frac{n!} 
{(n - 1)!}
\]}
{\(n\)}{\(\frac{n}
{n-1}\)}{\(\frac{n!}
{n!-1!}\)}{\(- 1\)}{}{}}
\LANGpl{
\Question{
Uprość dla \(n\in \mathbb{N}\).
\[ 
 \frac{n!} 
{(n - 1)!}
\]}
{\(n\)}{\(\frac{n}
{n-1}\)}{\(\frac{n!}
{n!-1!}\)}{\(- 1\)}{}{}}
\LANGcs{
\Question{
Z nabízených možností vyberte výraz, který se pro všechna
\(n\in \mathbb{N}\) rovná
výrazu \(\frac{n!}
{(n-1)!}\).}
{\(n\)}{\(\frac{n}
{n-1}\)}{\(\frac{n!}
{n!-1!}\)}{\(- 1\)}{}{}}
\LANGsk{
\Question{
Zjednodušte pre \(n\in \mathbb{N}\).
\[ 
 \frac{n!} 
{(n - 1)!}
\]}
{\(n\)}{\(\frac{n}
{n-1}\)}{\(\frac{n!}
{n!-1!}\)}{\(- 1\)}{}{}}
\LANGes{
\Question{
Simplifica esta expresión para que \(n\in \mathbb{N}\).
\[ 
 \frac{n!} 
{(n - 1)!}
\]}
{\(n\)}{\(\frac{n}
{n-1}\)}{\(\frac{n!}
{n!-1!}\)}{\(- 1\)}{}{}}
\End

\Data\ID{9000140502}
\Updated{2020-12-03 06:12:00}
\Level{B}
\SubArea{Combinatorics}
\LANGen{
\Question{
Simplify for \(n\in \mathbb{N}\).
\[ 
 \frac{(n + 1)!} 
{(n - 1)!}
\]}
{\(n^{2} + n\)}{\((n + 1)^{2}\)}{\(\frac{n+1}
{n-1}\)}{\(- 1\)}{}{}}
\LANGpl{
\Question{
Uprość dla \(n\in \mathbb{N}\).
\[ 
 \frac{(n + 1)!} 
{(n - 1)!}
\]}
{\(n^{2} + n\)}{\((n + 1)^{2}\)}{\(\frac{n+1}
{n-1}\)}{\(- 1\)}{}{}}
\LANGcs{
\Question{
Z nabízených možností vyberte výraz, který se pro všechna
\(n\in \mathbb{N}\) rovná
výrazu \(\frac{(n+1)!}
{(n-1)!}\).}
{\(n^{2} + n\)}{\((n + 1)^{2}\)}{\(\frac{n+1}
{n-1}\)}{\(- 1\)}{}{}}
\LANGsk{
\Question{
Zjednodušte pre \(n\in \mathbb{N}\).
\[ 
 \frac{(n + 1)!} 
{(n - 1)!}
\]}
{\(n^{2} + n\)}{\((n + 1)^{2}\)}{\(\frac{n+1}
{n-1}\)}{\(- 1\)}{}{}}
\LANGes{
\Question{
Simplifica esta expresión para que \(n\in \mathbb{N}\).
\[ 
 \frac{(n + 1)!} 
{(n - 1)!}
\]}
{\(n^{2} + n\)}{\((n + 1)^{2}\)}{\(\frac{n+1}
{n-1}\)}{\(- 1\)}{}{}}
\End

\Data\ID{9000140503}
\Updated{2020-12-03 06:12:01}
\Level{B}
\SubArea{Combinatorics}
\LANGen{
\Question{
Simplify for \(n\in \mathbb{N}\).
\[ 
 \frac{(n + 1)! + (n - 1)!} 
 {n!}
\]}
{\(\frac{n^{2}+n+1}
 {n} \)}{\(2\)}{\(n^{2} - 1\)}{\(\frac{n^{2}-n+1}
 {n} \)}{}{}}
\LANGpl{
\Question{
Uprość dla \(n\in \mathbb{N}\).
\[ 
 \frac{(n + 1)! + (n - 1)!} 
 {n!}
\]}
{\(\frac{n^{2}+n+1}
 {n} \)}{\(2\)}{\(n^{2} - 1\)}{\(\frac{n^{2}-n+1}
 {n} \)}{}{}}
\LANGcs{
\Question{
Z nabízených možností vyberte výraz, který se pro všechna
\(n\in \mathbb{N}\) rovná
výrazu \(\frac{(n+1)!+(n-1)!}
 {n!} \).}
{\(\frac{n^{2}+n+1}
 {n} \)}{\(2\)}{\(n^{2} - 1\)}{\(\frac{n^{2}-n+1}
 {n} \)}{}{}}
\LANGsk{
\Question{
Zjednodušte pre \(n\in \mathbb{N}\).
\[ 
 \frac{(n + 1)! + (n - 1)!} 
 {n!}
\]}
{\(\frac{n^{2}+n+1}
 {n} \)}{\(2\)}{\(n^{2} - 1\)}{\(\frac{n^{2}-n+1}
 {n} \)}{}{}}
\LANGes{
\Question{
Simplifica esta expresión para que \(n\in \mathbb{N}\).
\[ 
 \frac{(n + 1)! + (n - 1)!} 
 {n!}
\]}
{\(\frac{n^{2}+n+1}
 {n} \)}{\(2\)}{\(n^{2} - 1\)}{\(\frac{n^{2}-n+1}
 {n} \)}{}{}}
\End

\Data\ID{9000140504}
\Updated{2020-08-19 11:08:43}
\Level{B}
\SubArea{Combinatorics}
\LANGen{
\Question{
Simplify for \(n\in \mathbb{N}\),
\(n\geq 2\).
\[ 
 \frac{n\cdot (n - 2)!} 
{(n - 1)\cdot n!}
\]}
{\(\frac{1}
{(n-1)^{2}} \)}{\(\frac{(n^{2}-2n)!}
 {(n^{2}-n)!} \)}{\(\frac{n+1}
{n-1}\)}{\(\frac{(n-2)!}
{(n-1)!}\)}{}{}}
\LANGpl{
\Question{
Uprość dla \(n\in \mathbb{N}\),
\(n\geq 2\).
\[ 
 \frac{n\cdot (n - 2)!} 
{(n - 1)\cdot n!}
\]}
{\(\frac{1}
{(n-1)^{2}} \)}{\(\frac{(n^{2}-2n)!}
 {(n^{2}-n)!} \)}{\(\frac{n+1}
{n-1}\)}{\(\frac{(n-2)!}
{(n-1)!}\)}{}{}}
\LANGcs{
\Question{
Z nabízených možností vyberte výraz, který se pro všechna
\(n\in \mathbb{N}\), $n\geq2$ rovná
výrazu \(\frac{n\cdot (n-2)!}
{(n-1)\cdot n!}\).}
{\(\frac{1}
{(n-1)^{2}} \)}{\(\frac{(n^{2}-2n)!}
 {(n^{2}-n)!} \)}{\(\frac{n+1}
{n-1}\)}{\(\frac{(n-2)!}
{(n-1)!}\)}{}{}}
\LANGsk{
\Question{
Zjednodušte pre \(n\in \mathbb{N}\),
\(n\geq 2\).
\[ 
 \frac{n\cdot (n - 2)!} 
{(n - 1)\cdot n!}
\]}
{\(\frac{1}
{(n-1)^{2}} \)}{\(\frac{(n^{2}-2n)!}
 {(n^{2}-n)!} \)}{\(\frac{n+1}
{n-1}\)}{\(\frac{(n-2)!}
{(n-1)!}\)}{}{}}
\LANGes{
\Question{
Simplifica esta expresión para que \(n\in \mathbb{N}\),
\(n\geq 2\).
\[ 
 \frac{n\cdot (n - 2)!} 
{(n - 1)\cdot n!}
\]}
{\(\frac{1}
{(n-1)^{2}} \)}{\(\frac{(n^{2}-2n)!}
 {(n^{2}-n)!} \)}{\(\frac{n+1}
{n-1}\)}{\(\frac{(n-2)!}
{(n-1)!}\)}{}{}}
\End

\Data\ID{9000140505}
\Updated{2020-08-19 11:08:59}
\Level{B}
\SubArea{Combinatorics}
\LANGen{
\Question{
Simplify \(\frac{72!}
{70!+71!}\).}
{\(71\)}{\(72\)}{\(\frac{72}
{141}\)}{\(\frac{72!}
{141!}\)}{}{}}
\LANGpl{
\Question{
Uprość \(\frac{72!}
{70!+71!}\).}
{\(71\)}{\(72\)}{\(\frac{72}
{141}\)}{\(\frac{72!}
{141!}\)}{}{}}
\LANGcs{
\Question{
Z nabízených možností vyberte výraz, který se rovná výrazu
\(\frac{72!}
{70!+71!}\).}
{\(71\)}{\(72\)}{\(\frac{72}
{141}\)}{\(\frac{72!}
{141!}\)}{}{}}
\LANGsk{
\Question{
Zjednodušte \(\frac{72!}
{70!+71!}\).}
{\(71\)}{\(72\)}{\(\frac{72}
{141}\)}{\(\frac{72!}
{141!}\)}{}{}}
\LANGes{
\Question{
Simplifica \(\frac{72!}
{70!+71!}\).}
{\(71\)}{\(72\)}{\(\frac{72}
{141}\)}{\(\frac{72!}
{141!}\)}{}{}}
\End

\Data\ID{9000140506}
\Updated{2020-08-19 11:08:15}
\Level{B}
\SubArea{Combinatorics}
\LANGen{
\Question{
Simplify \(\frac{2!}
{1!} + \frac{3!} 
{2!} + \frac{4!} 
{3!}\).}
{\(9\)}{\(\frac{29}
 {6} \)}{\(\frac{9}
{6}\)}{\(\frac{12!+9!+8!}
 {6!} \)}{}{}}
\LANGpl{
\Question{
Uprość  \(\frac{2!}
{1!} + \frac{3!} 
{2!} + \frac{4!} 
{3!}\).}
{\(9\)}{\(\frac{29}
 {6} \)}{\(\frac{9}
{6}\)}{\(\frac{12!+9!+8!}
 {6!} \)}{}{}}
\LANGcs{
\Question{
Z nabízených možností vyberte výraz, který se rovná výrazu
\(\frac{2!}
{1!} + \frac{3!} 
{2!} + \frac{4!} 
{3!}\).}
{\(9\)}{\(\frac{29}
 {6} \)}{\(\frac{9}
{6}\)}{\(\frac{12!+9!+8!}
 {6!} \)}{}{}}
\LANGsk{
\Question{
Zjednodušte \(\frac{2!}
{1!} + \frac{3!} 
{2!} + \frac{4!} 
{3!}\).}
{\(9\)}{\(\frac{29}
 {6} \)}{\(\frac{9}
{6}\)}{\(\frac{12!+9!+8!}
 {6!} \)}{}{}}
\LANGes{
\Question{
Simplifica \(\frac{2!}
{1!} + \frac{3!} 
{2!} + \frac{4!} 
{3!}\).}
{\(9\)}{\(\frac{29}
 {6} \)}{\(\frac{9}
{6}\)}{\(\frac{12!+9!+8!}
 {6!} \)}{}{}}
\End

\Data\ID{9000140507}
\Updated{2020-08-19 11:08:33}
\Level{B}
\SubArea{Combinatorics}
\LANGen{
\Question{
Simplify for \(n\in \mathbb{N}\).
\[ 
 \frac{(n + 2)!} 
{n! + (n + 1)!}
\]}
{\(n + 1\)}{\(\frac{n!+2}
{2n!+1}\)}{\(\frac{n+2}
{2n+1}\)}{\(\frac{n!+2!}
{2n!+1!}\)}{}{}}
\LANGpl{
\Question{
Uprość dla \(n\in \mathbb{N}\).
\[ 
 \frac{(n + 2)!} 
{n! + (n + 1)!}
\]}
{\(n + 1\)}{\(\frac{n!+2}
{2n!+1}\)}{\(\frac{n+2}
{2n+1}\)}{\(\frac{n!+2!}
{2n!+1!}\)}{}{}}
\LANGcs{
\Question{
Z nabízených možností vyberte výraz, který se pro všechna
\(n\in \mathbb{N}\) rovná
výrazu \(\frac{(n+2)!}
{n!+(n+1)!}\).}
{\(n + 1\)}{\(\frac{n!+2}
{2n!+1}\)}{\(\frac{n+2}
{2n+1}\)}{\(\frac{n!+2!}
{2n!+1!}\)}{}{}}
\LANGsk{
\Question{
Zjednodušte pre \(n\in \mathbb{N}\).
\[ 
 \frac{(n + 2)!} 
{n! + (n + 1)!}
\]}
{\(n + 1\)}{\(\frac{n!+2}
{2n!+1}\)}{\(\frac{n+2}
{2n+1}\)}{\(\frac{n!+2!}
{2n!+1!}\)}{}{}}
\LANGes{
\Question{
Simplifica esta expresión para que   \(n\in \mathbb{N}\).
\[ 
 \frac{(n + 2)!} 
{n! + (n + 1)!}
\]}
{\(n + 1\)}{\(\frac{n!+2}
{2n!+1}\)}{\(\frac{n+2}
{2n+1}\)}{\(\frac{n!+2!}
{2n!+1!}\)}{}{}}
\End

\Data\ID{9000140508}
\Updated{2020-08-19 11:08:45}
\Level{B}
\SubArea{Combinatorics}
\LANGen{
\Question{
Simplify for \(n\in \mathbb{N}\).
\[ 
 \frac{(n + 1)!} 
{n! - (n + 1)!}
\]}
{\(- 1 - \frac{1} 
{n}\)}{\(n + 1\)}{\(n! + 1\)}{\(- \frac{n+1} 
{(n-1)!}\)}{}{}}
\LANGpl{
\Question{
Uprość dla \(n\in \mathbb{N}\).
\[ 
 \frac{(n + 1)!} 
{n! - (n + 1)!}
\]}
{\(- 1 - \frac{1} 
{n}\)}{\(n + 1\)}{\(n! + 1\)}{\(- \frac{n+1} 
{(n-1)!}\)}{}{}}
\LANGcs{
\Question{
Z nabízených možností vyberte výraz, který se pro všechna
\(n\in \mathbb{N}\) rovná
výrazu \(\frac{(n+1)!}
{n!-(n+1)!}\).}
{\(- 1 - \frac{1} 
{n}\)}{\(n + 1\)}{\(n! + 1\)}{\(- \frac{n+1} 
{(n-1)!}\)}{}{}}
\LANGsk{
\Question{
Zjednodušte pre \(n\in \mathbb{N}\).
\[ 
 \frac{(n + 1)!} 
{n! - (n + 1)!}
\]}
{\(- 1 - \frac{1} 
{n}\)}{\(n + 1\)}{\(n! + 1\)}{\(- \frac{n+1} 
{(n-1)!}\)}{}{}}
\LANGes{
\Question{
Simplifica esta expresión para que   \(n\in \mathbb{N}\).
\[ 
 \frac{(n + 1)!} 
{n! - (n + 1)!}
\]}
{\(- 1 - \frac{1} 
{n}\)}{\(n + 1\)}{\(n! + 1\)}{\(- \frac{n+1} 
{(n-1)!}\)}{}{}}
\End

\Data\ID{9000140509}
\Updated{2021-04-25 07:04:35}
\Level{B}
\SubArea{Combinatorics}
\LANGen{
\Question{
Find the solution set of the equation assuming
\(x\in \mathbb{N}\).
\[ 
 (x + 1)! = 6(x - 1)!
\]}
{\(\{2\}\)}{\(\{ - 3;\ 2\}\)}{\(\left \{\frac{5}
{7}\right \}\)}{\(\left \{\frac{7}
{5}\right \}\)}{}{}}
\LANGpl{
\Question{
Wyznacz zbiór rozwiązań równania dla 
\(x\in \mathbb{N}\).
\[ 
 (x + 1)! = 6(x - 1)!
\]}
{\(\{2\}\)}{\(\{ - 3;\ 2\}\)}{\(\left \{\frac{5}
{7}\right \}\)}{\(\left \{\frac{7}
{5}\right \}\)}{}{}}
\LANGcs{
\Question{
Určete množinu kořenů dané rovnice pro 
\(x\in \mathbb{N}\).
\[ 
 (x + 1)! = 6(x - 1)!
\]}
{\(\{2\}\)}{\(\{ - 3;\ 2\}\)}{\(\left \{\frac{5}
{7}\right \}\)}{\(\left \{\frac{7}
{5}\right \}\)}{}{}}
\LANGsk{
\Question{
Nájdite množinu koreňov danej rovnice pre
\(x\in \mathbb{N}\).
\[ 
 (x + 1)! = 6(x - 1)!
\]}
{\(\{2\}\)}{\(\{ - 3;\ 2\}\)}{\(\left \{\frac{5}
{7}\right \}\)}{\(\left \{\frac{7}
{5}\right \}\)}{}{}}
\LANGes{
\Question{
Halla el conjunto de soluciones (o solución) de la ecuación suponiendo que
\(x\in \mathbb{N}\).
\[ 
 (x + 1)! = 6(x - 1)!
\]}
{\(\{2\}\)}{\(\{ - 3;\ 2\}\)}{\(\left \{\frac{5}
{7}\right \}\)}{\(\left \{\frac{7}
{5}\right \}\)}{}{}}
\End

\Data\ID{9000140510}
\Updated{2021-04-25 07:04:22}
\Level{B}
\SubArea{Combinatorics}
\LANGen{
\Question{
Find the solution set of the equation assuming
\(x\in \mathbb{N}\).
\[ 
 (x + 1)! + x! = 6x + 12
\]}
{\(\{3\}\)}{\(\{2\}\)}{\(\{1\}\)}{\(\{4\}\)}{}{}}
\LANGpl{
\Question{
Wyznacz zbiór rozwiązań równania dla
\(x\in \mathbb{N}\).
\[ 
 (x + 1)! + x! = 6x + 12
\]}
{\(\{3\}\)}{\(\{2\}\)}{\(\{1\}\)}{\(\{4\}\)}{}{}}
\LANGcs{
\Question{
Určete množinu kořenů dané rovnice pro 
\(x\in \mathbb{N}\).
\[ 
 (x + 1)! + x! = 6x + 12
\]}
{\(\{3\}\)}{\(\{2\}\)}{\(\{1\}\)}{\(\{4\}\)}{}{}}
\LANGsk{
\Question{
Nájdite množinu koreňov danej rovnice pre
\(x\in \mathbb{N}\).
\[ 
 (x + 1)! + x! = 6x + 12
\]}
{\(\{3\}\)}{\(\{2\}\)}{\(\{1\}\)}{\(\{4\}\)}{}{}}
\LANGes{
\Question{
Halla la solución de la ecuación suponiendo que  
\(x\in \mathbb{N}\).
\[ 
 (x + 1)! + x! = 6x + 12
\]}
{\(\{3\}\)}{\(\{2\}\)}{\(\{1\}\)}{\(\{4\}\)}{}{}}
\End

\Data\ID{9000141501}
\Updated{2021-04-25 07:04:24}
\Level{B}
\SubArea{Combinatorics}
\LANGen{
\Question{
Let \(A\) be set with
\(n\) mutually different
elements. If \(n\) is increased
by \(2\), then number
of \(3\)-permutations
is increased by \(384\).
Find \(n\). (The term
„\(k\)-permutation” stands for
an ordered arrangement of \(k\) 
objects from a set of \(n\) 
objects.)}
{\(8\)}{\(64\)}{\(32\)}{}{}{}}
\LANGpl{
\Question{
\(A\) to zbiór 
\(n\) różnych elementów. Jeśli \(n\) wzrośnie o 
 \(2\), wtedy liczba 
 \(3\)-elementów wariacji zbioru  \(n\)  wzrośnie o
 \(384\).
Wyznacz  \(n\).}
{\(8\)}{\(64\)}{\(32\)}{}{}{}}
\LANGcs{
\Question{
Zvětší-li se počet prvků o \(2\),
zvětší se počet z nich vytvořených variací
\(3\). třídy bez
opakování o \(384\).
Určete původní počet prvků.}
{\(8\)}{\(64\)}{\(32\)}{}{}{}}
\LANGsk{
\Question{
Nech \(A\) je množina s 
\(n\)  navzájom rozdielnymi prvkami. Ak sa počet prvkov \(n\) zväčší o \(2\), potom počet z nich vytvorených variácií
 \(3\). triedy bez opakovania sa zväčší o
 \(384\).
Určte pôvodný počet prvkov \(n\).}
{\(8\)}{\(64\)}{\(32\)}{}{}{}}
\LANGes{
\Question{
Sea \(A\) un conjunto con
\(n\) elementos diferentes. Si \(n\) aumenta en \(2\), luego el número de permutaciones de \(3\) elementos aumenta en 
\(384\).
Calcula \(n\). (El término 
„\(k\)-permutación” significa
una disposición ordenada de \(k\) 
objetos de un conjunto de \(n\) 
objetos.)}
{\(8\)}{\(64\)}{\(32\)}{}{}{}}
\End

\Data\ID{9000141502}
\Updated{2021-04-25 07:04:01}
\Level{B}
\SubArea{Combinatorics}
\LANGen{
\Question{
Let \(A\) be set with
\(n\) mutually different elements.
The number of \(5\)-permutations
with repetition is \(1024\).
Find \(n\). (The term
„\(k\)-permutation
with repetition” stands for an ordered arrangement of
\(k\) objects from
a set of \(n\) 
objects, when each object can be chosen more than once.)}
{\(4\)}{\(5\)}{\(2\)}{}{}{}}
\LANGpl{
\Question{
\(A\) to zbiór różnych elementów
\(n\). 
Liczba   \(5\) elementówych  wariacji z powtórzeniami  \(A\)  \(n\)-elementówego wynosi
wynosi \(1024\).
Wyznacz \(n\).}
{\(4\)}{\(5\)}{\(2\)}{}{}{}}
\LANGcs{
\Question{
Z kolika prvků lze vytvořit \(1024\) 
variací \(5\).
třídy s opakováním?}
{\(4\)}{\(5\)}{\(2\)}{}{}{}}
\LANGsk{
\Question{
Nech \(A\) je množina s
\(n\) navzájom rozdielnymi prvkami.
Počet variácií \(5\). triedy s opakovaním je \(1024\).
Určte počet prvkov \(n\).}
{\(4\)}{\(5\)}{\(2\)}{}{}{}}
\LANGes{
\Question{
Sea \(A\) un conjunto con \(n\) elementos diferentes.
El número de permutaciones de \(5\) elementos
con repetición es \(1024\).
Calcula \(n\). (El término
„\(k\)-permutación con repetición” significa una disposición ordenada de \(k\)  objetos de un conjunto de \(n\)  objetos, donde cada objeto puede ser elegido más de una vez).}
{\(4\)}{\(5\)}{\(2\)}{}{}{}}
\End

\Data\ID{9000141503}
\Updated{2021-04-25 07:04:02}
\Level{B}
\SubArea{Combinatorics}
\LANGen{
\Question{
Let \(A\) be a
set with \(n\) 
mutually different elements. If two elements are removed from
\(A\), the number of all
permutations of set \(A\) 
decreases \(20\)-times.
Find \(n\).}
{\(5\)}{\(4\)}{\(5\) or
\(- 4\)}{}{}{}}
\LANGpl{
\Question{
\(A\) to zbiór różnych elementów  \(n\). Jeśli usuniemy dwa elementy ze zbioru 
\(A\), to liczba wszystkich permutacji zbioru \(A\) 
zmniejszy się \(20\)-krotnie.
Wyznacz \(n\).}
{\(5\)}{\(4\)}{\(5\) lub
\(- 4\)}{}{}{}}
\LANGcs{
\Question{
Zmenší-li se počet prvků o \(2\),
zmenší se počet z nich vytvořených permutací bez opakování
dvacetkrát. Určete původní počet prvků.}
{\(5\)}{\(4\)}{$5$ nebo $ - 4$}{}{}{}}
\LANGsk{
\Question{
Nech \(A\) je množina s \(n\) navzájom rôznymi prvkami. Ak sa zmenší počet prvkov množiny
\(A\) o dva, zmenší sa počet z nich vytvorených permutácií množiny
 \(A\) 
 \(20\)-krát.
Určte pôvodný počet prvkov \(n\).}
{\(5\)}{\(4\)}{\(5\) alebo
\(- 4\)}{}{}{}}
\LANGes{
\Question{
Sea \(A\) un conjunto con \(n\) elementos diferentes. Si dos elementos son eliminados del conjunto 
\(A\), el número de todas las permutaciones del conjunto \(A\) 
disminuye en  \(20\) .
Calcula \(n\).}
{\(5\)}{\(4\)}{\(5\) or
\(- 4\)}{}{}{}}
\End

\Data\ID{9000141504}
\Updated{2021-04-25 07:04:45}
\Level{B}
\SubArea{Combinatorics}
\LANGen{
\Question{
Let \(A\) be a
set with \(n\) 
mutually different elements. If we add one element to the set
\(A\), the number of
\(3\)-combinations
of a set \(A\) is
increased by \(21\).
Find \(n\).}
{\(7\)}{\(6\)}{\(43\)}{}{}{}}
\LANGpl{
\Question{
\(A\) to zbiór różnych elementów
 \(n\). Jeśli dodamy jeden element do zbioru
\(A\), liczba 
\(3\)-elementówych kombinacji 
zbioru \(A\) zostanie
zwiększona o \(21\).
Wyznacz \(n\).}
{\(7\)}{\(6\)}{\(43\)}{}{}{}}
\LANGcs{
\Question{
Zvětší-li se počet prvků o \(1\),
zvětší se počet z nich vytvořených kombinací
\(3\). třídy bez
opakování o \(21\).
Určete původní počet prvků.}
{\(7\)}{\(6\)}{\(43\)}{}{}{}}
\LANGsk{
\Question{
Nech \(A\) je množina, ktorá má \(n\) navzájom rôznych prvkov. Ak pridáme jeden prvok k množine
\(A\), počet kombinácií
\(3\). triedy
množiny \(A\) sa zväčší
o \(21\).
Určte pôvodný počet prvkov \(n\).}
{\(7\)}{\(6\)}{\(43\)}{}{}{}}
\LANGes{
\Question{
Sea \(A\) un conjunto con \(n\)  elementos diferentes. Si  añadimos un elemento al conjunto  
\(A\), el número de 
\(3\) combinaciones
del conjunto \(A\) aumenta en \(21\).
Calcula \(n\).}
{\(7\)}{\(6\)}{\(43\)}{}{}{}}
\End

\Data\ID{9000141505}
\Updated{2020-08-19 11:08:39}
\Level{B}
\SubArea{Combinatorics}
\LANGen{
\Question{
Assuming \(x\in \mathbb{N}\),
find the solution set of the following equation. 
\[ 
 \frac{(x + 2)!} 
 {x!} = 2\cdot \frac{x!}
{(x - 2)!} + 3!
\]}
{\(\{4\}\)}{\(\{4;1\}\)}{\(\{1\}\)}{}{}{}}
\LANGpl{
\Question{
Dla \(x\in \mathbb{N}\),
wyznacz zbiór rozwiązań podanego równania. 
\[ 
 \frac{(x + 2)!} 
 {x!} = 2\cdot \frac{x!}
{(x - 2)!} + 3!
\]}
{\(\{4\}\)}{\(\{4;1\}\)}{\(\{1\}\)}{}{}{}}
\LANGcs{
\Question{
Určete množinu všech řešení následující rovnice.
\[ 
\frac{(x+2)!}
 {x!} = 2\cdot \frac{x!}
{(x-2)!} + 3!
\]}
{\(\{4\}\)}{\(\{4;1\}\)}{\(\{1\}\)}{}{}{}}
\LANGsk{
\Question{
Nech \(x\in \mathbb{N}\).
Určte množinu všetkých riešení danej rovnice. 
\[ 
 \frac{(x + 2)!} 
 {x!} = 2\cdot \frac{x!}
{(x - 2)!} + 3!
\]}
{\(\{4\}\)}{\(\{4;1\}\)}{\(\{1\}\)}{}{}{}}
\LANGes{
\Question{
Suponiendo que  \(x\in \mathbb{N}\),
halla la solución de la siguiente ecuación. 
\[ 
 \frac{(x + 2)!} 
 {x!} = 2\cdot \frac{x!}
{(x - 2)!} + 3!
\]}
{\(\{4\}\)}{\(\{4;1\}\)}{\(\{1\}\)}{}{}{}}
\End

\Data\ID{9000141506}
\Updated{2020-08-19 12:08:39}
\Level{B}
\SubArea{Combinatorics}
\LANGen{
\Question{
Assuming \(x\in \mathbb{N}\),
find the solution set of the following equation. 
\[ 
 \left({11\above 0.0pt 
 4} \right) =\left ({11\above 0.0pt 
 x} \right)
\]}
{\(\{4;7\}\)}{\(\{4\}\)}{\(\{ - 4\}\)}{}{}{}}
\LANGpl{
\Question{
Dla \(x\in \mathbb{N}\),
wyznacz zbiór rozwiązań podanego równania. 
\[ 
 \left({11\above 0.0pt 
 4} \right) =\left ({11\above 0.0pt 
 x} \right)
\]}
{\(\{4;7\}\)}{\(\{4\}\)}{\(\{ - 4\}\)}{}{}{}}
\LANGcs{
\Question{
Určete množinu všech řešení následující rovnice.
\[ \left({11\above 0.0pt 
 4} \right) =\left ({11\above 0.0pt 
 x} \right)\]}
{\(\{4;7\}\)}{\(\{4\}\)}{\(\{ - 4\}\)}{}{}{}}
\LANGsk{
\Question{
Nech \(x\in \mathbb{N}\).
Určte množinu všetkých riešení danej rovnice. 
\[ 
 \left({11\above 0.0pt 
 4} \right) =\left ({11\above 0.0pt 
 x} \right)
\]}
{\(\{4;7\}\)}{\(\{4\}\)}{\(\{ - 4\}\)}{}{}{}}
\LANGes{
\Question{
Suponiendo que \(x\in \mathbb{N}\),
halla la solución de la siguiente ecuación.
\[ 
 \left({11\above 0.0pt 
 4} \right) =\left ({11\above 0.0pt 
 x} \right)
\]}
{\(\{4;7\}\)}{\(\{4\}\)}{\(\{ - 4\}\)}{}{}{}}
\End

\Data\ID{9000141507}
\Updated{2020-08-19 12:08:55}
\Level{B}
\SubArea{Combinatorics}
\LANGen{
\Question{
Assuming \(x,y\in \mathbb{N}\),
find the solution set of the following equation. 
\[ 
 \left({x\above 0.0pt 
y}\right)^{2} - 2\cdot \left({x\above 0.0pt 
y}\right) - 3 = 0
\]}
{\(\{[3;1];[3;2]\}\)}{\(\{[3;1]\}\)}{\(\{[3;1];[1;3]\}\)}{}{}{}}
\LANGpl{
\Question{
Dla \(x,y\in \mathbb{N}\),
wyznacz zbiór rozwiązań podanego równania. 
\[ 
 \left({x\above 0.0pt 
y}\right)^{2} - 2\cdot \left({x\above 0.0pt 
y}\right) - 3 = 0
\]}
{\(\{[3;1];[3;2]\}\)}{\(\{[3;1]\}\)}{\(\{[3;1];[1;3]\}\)}{}{}{}}
\LANGcs{
\Question{
Určete množinu všech řešení dané rovnice.
\[ \left({x\above 0.0pt 
y}\right)^{2} - 2\cdot \left({x\above 0.0pt 
y}\right) - 3 = 0\]}
{\(\{[3;1];[3;2]\}\)}{\(\{[3;1]\}\)}{\(\{[3;1];[1;3]\}\)}{}{}{}}
\LANGsk{
\Question{
Nech \(x,y\in \mathbb{N}\).
Určte množinu všetkých riešení danej rovnice. 
\[ 
 \left({x\above 0.0pt 
y}\right)^{2} - 2\cdot \left({x\above 0.0pt 
y}\right) - 3 = 0
\]}
{\(\{[3;1];[3;2]\}\)}{\(\{[3;1]\}\)}{\(\{[3;1];[1;3]\}\)}{}{}{}}
\LANGes{
\Question{
Suponiendo que \(x,y\in \mathbb{N}\),
halla la solución de la siguiente ecuación.
\[ 
 \left({x\above 0.0pt 
y}\right)^{2} - 2\cdot \left({x\above 0.0pt 
y}\right) - 3 = 0
\]}
{\(\{[3;1];[3;2]\}\)}{\(\{[3;1]\}\)}{\(\{[3;1];[1;3]\}\)}{}{}{}}
\End

\Data\ID{9000141508}
\Updated{2020-08-19 12:08:05}
\Level{B}
\SubArea{Combinatorics}
\LANGen{
\Question{
Assuming \(x\in \mathbb{N}\),
find the solution set of the following equation. 
\[ 
 \left({x\above 0.0pt 
x}\right) +\left ({x + 1\above 0.0pt 
 x} \right) +\left ({x + 2\above 0.0pt 
 x} \right) +\left ({x + 3\above 0.0pt 
 x} \right) = \frac{x^{3} + 59} 
 {6}
\]}
{\(\{1\}\)}{\(\{4\}\)}{\(\{10\}\)}{}{}{}}
\LANGpl{
\Question{
Dla \(x\in \mathbb{N}\), wyznacz zbiór rozwiązań podanego równania. 
\[ 
 \left({x\above 0.0pt 
x}\right) +\left ({x + 1\above 0.0pt 
 x} \right) +\left ({x + 2\above 0.0pt 
 x} \right) +\left ({x + 3\above 0.0pt 
 x} \right) = \frac{x^{3} + 59} 
 {6}
\]}
{\(\{1\}\)}{\(\{4\}\)}{\(\{10\}\)}{}{}{}}
\LANGcs{
\Question{
Určete množinu všech řešení dané rovnice.
\[ \left({x\above 0.0pt 
x}\right) +\left ({x+1\above 0.0pt 
 x} \right) +\left ({x+2\above 0.0pt 
 x} \right) +\left ({x+3\above 0.0pt 
 x} \right) = \frac{x^{3}+59} 
 {6} \]}
{\(\{1\}\)}{\(\{4\}\)}{\(\{10\}\)}{}{}{}}
\LANGsk{
\Question{
Nech \(x\in \mathbb{N}\).
Určte množinu všetkých riešení danej rovnice. 
\[ 
 \left({x\above 0.0pt 
x}\right) +\left ({x + 1\above 0.0pt 
 x} \right) +\left ({x + 2\above 0.0pt 
 x} \right) +\left ({x + 3\above 0.0pt 
 x} \right) = \frac{x^{3} + 59} 
 {6}
\]}
{\(\{1\}\)}{\(\{4\}\)}{\(\{10\}\)}{}{}{}}
\LANGes{
\Question{
Suponiendo que \(x\in \mathbb{N}\),
halla la solución de la siguiente ecuación.
\[ 
 \left({x\above 0.0pt 
x}\right) +\left ({x + 1\above 0.0pt 
 x} \right) +\left ({x + 2\above 0.0pt 
 x} \right) +\left ({x + 3\above 0.0pt 
 x} \right) = \frac{x^{3} + 59} 
 {6}
\]}
{\(\{1\}\)}{\(\{4\}\)}{\(\{10\}\)}{}{}{}}
\End

\Data\ID{9000141509}
\Updated{2020-08-19 12:08:21}
\Level{B}
\SubArea{Combinatorics}
\LANGen{
\Question{
Assuming \(x\in \mathbb{N}\),
find the solution set of the following inequality. 
\[ 
 2\cdot \left({x - 1\above 0.0pt 
x - 3}\right) + x\cdot (x - 9)\leq - 8
\]}
{\(\{3;4;5\}\)}{\(\{1;2;3;4;5\}\)}{\([ 1;5] \)}{}{}{}}
\LANGpl{
\Question{
Dla \(x\in \mathbb{N}\),
wyznacz zbiór rozwiązań podanej nierówności.
\[ 
 2\cdot \left({x - 1\above 0.0pt 
x - 3}\right) + x\cdot (x - 9)\leq - 8
\]}
{\(\{3;4;5\}\)}{\(\{1;2;3;4;5\}\)}{\([ 1;5] \)}{}{}{}}
\LANGcs{
\Question{
Určete množinu všech řešení dané nerovnice.
\[ 2\cdot \left({x-1\above 0.0pt 
x-3}\right) + x\cdot (x - 9)\leq - 8\]}
{\(\{3;4;5\}\)}{\(\{1;2;3;4;5\}\)}{\(\langle 1;5\rangle \)}{}{}{}}
\LANGsk{
\Question{
Nech \(x\in \mathbb{N}\).
Určte množinu všetkých riešení danej nerovnice. 
\[ 
 2\cdot \left({x - 1\above 0.0pt 
x - 3}\right) + x\cdot (x - 9)\leq - 8
\]}
{\(\{3;4;5\}\)}{\(\{1;2;3;4;5\}\)}{\([ 1;5] \)}{}{}{}}
\LANGes{
\Question{
Suponiendo que \(x\in \mathbb{N}\),
halla la solución de la siguiente inecuación.
\[ 
 2\cdot \left({x - 1\above 0.0pt 
x - 3}\right) + x\cdot (x - 9)\leq - 8
\]}
{\(\{3;4;5\}\)}{\(\{1;2;3;4;5\}\)}{\([ 1;5] \)}{}{}{}}
\End

\Data\ID{9000141510}
\Updated{2020-08-19 12:08:56}
\Level{B}
\SubArea{Combinatorics}
\LANGen{
\Question{
Assuming \(x\in \mathbb{N}\),
find the solution set of the following inequality. 
\[ 
 \left({ x\above 0.0pt 
x - 2}\right)\cdot \left({x\above 0.0pt 
2}\right) - 11\cdot \left({x\above 0.0pt 
2}\right) + 28 < 0
\]}
{\(\{4\}\)}{\(\{5;6\}\)}{\((4;7)\)}{}{}{}}
\LANGpl{
\Question{
Dla \(x\in \mathbb{N}\),
wyznacz zbiór rozwiązań podanej nierówności. 
\[ 
 \left({ x\above 0.0pt 
x - 2}\right)\cdot \left({x\above 0.0pt 
2}\right) - 11\cdot \left({x\above 0.0pt 
2}\right) + 28 < 0
\]}
{\(\{4\}\)}{\(\{5;6\}\)}{\((4;7)\)}{}{}{}}
\LANGcs{
\Question{
Určete množinu všech řešení dané nerovnice.
\[ \left({ x\above 0.0pt 
x-2}\right)\cdot \left({x\above 0.0pt 
2} \right) - 11\cdot \left({x\above 0.0pt 
2} \right) + 28 < 0\]}
{\(\{4\}\)}{\(\{5;6\}\)}{\((4;7)\)}{}{}{}}
\LANGsk{
\Question{
Nech \(x\in \mathbb{N}\).
Určte množinu všetkých riešení danej nerovnice. 
\[ 
 \left({ x\above 0.0pt 
x - 2}\right)\cdot \left({x\above 0.0pt 
2}\right) - 11\cdot \left({x\above 0.0pt 
2}\right) + 28 < 0
\]}
{\(\{4\}\)}{\(\{5;6\}\)}{\((4;7)\)}{}{}{}}
\LANGes{
\Question{
Suponiendo que \(x\in \mathbb{N}\), halla la solución de la siguiente inecuación.
\[ 
 \left({ x\above 0.0pt 
x - 2}\right)\cdot \left({x\above 0.0pt 
2}\right) - 11\cdot \left({x\above 0.0pt 
2}\right) + 28 < 0
\]}
{\(\{4\}\)}{\(\{5;6\}\)}{\((4;7)\)}{}{}{}}
\End

\Data\ID{1003024701}
\Updated{2020-07-19 04:07:06}
\Level{C}
\SubArea{Combinatorics}
\LANGen{
\Question{
What is the coefficient at \( x^7 \) in the expansion of \( (1-3x)^9 \)?}
{\( -78\:732 \)}{\( 59\:049 \)}{\( -59\:049 \)}{\( 78\:732 \)}{}{}}
\LANGpl{
\Question{
Jaki jest współczynnik w  \( x^7 \) w rozszerzeniu \( (1-3x)^9 \)?}
{\( -78\:732 \)}{\( 59\:049 \)}{\( -59\:049 \)}{\( 78\:732 \)}{}{}}
\LANGcs{
\Question{
Je daný výraz \( (1-3x)^9 \). Určete koeficient členu binomického rozvoje daného výrazu, který obsahuje \( x^7 \).}
{\( -78\:732 \)}{\( 59\:049 \)}{\( -59\:049 \)}{\( 78\:732 \)}{}{}}
\LANGsk{
\Question{
Je daný výraz \( (1-3x)^9 \). Určte koeficient toho člena binomického rozvoja daného výrazu, ktorý obsahuje \( x^7 \).}
{\( -78\:732 \)}{\( 59\:049 \)}{\( -59\:049 \)}{\( 78\:732 \)}{}{}}
\LANGes{
\Question{
¿Cuál es el coeficiente de \( x^7 \) en la expansión binomial de \( (1-3x)^9 \)?}
{\( -78\:732 \)}{\( 59\:049 \)}{\( -59\:049 \)}{\( 78\:732 \)}{}{}}
\End

\Data\ID{1003024702}
\Updated{2020-07-19 07:07:48}
\Level{C}
\SubArea{Combinatorics}
\LANGen{
\Question{
Find the constant term in the expansion of \( \left(2+3x^2\right)^{30} \).}
{\( 2^{30} \)}{\( 2^{29} \)}{\( 3\cdot2^{30} \)}{\( 3\cdot2^{29} \)}{}{}}
\LANGpl{
\Question{
Wyznacz element, który nie zawiera zmiennej  \(x\) w wyrażeniu  \( \left(2+3x^2\right)^{30} \).}
{\( 2^{30} \)}{\( 2^{29} \)}{\( 3\cdot2^{30} \)}{\( 3\cdot2^{29} \)}{}{}}
\LANGcs{
\Question{
V rozvoji výrazu \( \left(2+3x^2\right)^{30} \) určete člen, který neobsahuje proměnou \(x\).}
{\( 2^{30} \)}{\( 2^{29} \)}{\( 3\cdot2^{30} \)}{\( 3\cdot2^{29} \)}{}{}}
\LANGsk{
\Question{
V rozvoji výrazu \( \left(2+3x^2\right)^{30} \) určite člen, ktorý neobsahuje premennú \(x\).}
{\( 2^{30} \)}{\( 2^{29} \)}{\( 3\cdot2^{30} \)}{\( 3\cdot2^{29} \)}{}{}}
\LANGes{
\Question{
Determina el término independiente en la expansión de \( \left(2+3x^2\right)^{30} \).}
{\( 2^{30} \)}{\( 2^{29} \)}{\( 3\cdot2^{30} \)}{\( 3\cdot2^{29} \)}{}{}}
\End

\Data\ID{1003024703}
\Updated{2021-04-25 06:04:21}
\Level{C}
\SubArea{Combinatorics}
\LANGen{
\Question{
Let the term with \( x^{-3} \) in the expansion of \( \left(4x^{-\frac12}-\frac12\right)^{10} \) equals \( 105 \). Find the value of \( x \).}
{\( 8 \)}{\( 9 \)}{\( 10 \)}{\( 11 \)}{}{}}
\LANGpl{
\Question{
Załóżmy, że współczynnik w \( x^{-3} \) w rozszerzeniu \( \left(4x^{-\frac12}-\frac12\right)^{10} \) wynosi \( 105 \). Oblicz wartość \( x \).}
{\( 8 \)}{\( 9 \)}{\( 10 \)}{\( 11 \)}{}{}}
\LANGcs{
\Question{
Určete hodnotu proměnné \(x\), pro kterou se člen binomického rozvoje výrazu \( \left(4x^{-\frac12}-\frac12\right)^{10} \) obsahující \( x^{-3} \) rovná  \( 105 \).}
{\( 8 \)}{\( 9 \)}{\( 10 \)}{\( 11 \)}{}{}}
\LANGsk{
\Question{
Určite hodnotu premennej \(x\), pre ktorú sa člen binomického rozvoja výrazu \( \left(4x^{-\frac12}-\frac12\right)^{10} \) obsahujúci \( x^{-3} \) rovná  \( 105 \).}
{\( 8 \)}{\( 9 \)}{\( 10 \)}{\( 11 \)}{}{}}
\LANGes{
\Question{
Sabiendo que el término con \( x^{-3} \) en la expansión de \( \left(4x^{-\frac12}-\frac12\right)^{10} \) es igual a \( 105 \). Determina el valor de \( x \).}
{\( 8 \)}{\( 9 \)}{\( 10 \)}{\( 11 \)}{}{}}
\End

\Data\ID{1003024704}
\Updated{2020-07-19 07:07:40}
\Level{C}
\SubArea{Combinatorics}
\LANGen{
\Question{
Suppose the coefficient of the quadratic term (i.e. the coefficient at \( x^2 \)) in the expansion of \( (1+x)^n \) equals \( 300 \). Find the value of the exponent \( n \).}
{\( 25 \)}{\( 24 \)}{\( 30 \)}{\( 20 \)}{}{}}
\LANGpl{
\Question{
Załóżmy, że współczynnik wyrażenia kwadratowego (tj. współczynnik w \( x^2 \)) w rozszerzeniu  \( (1+x)^n \) wynosi \( 300 \). Oblicz wartość wykładnika \( n \).}
{\( 25 \)}{\( 24 \)}{\( 30 \)}{\( 20 \)}{}{}}
\LANGcs{
\Question{
Jaké je třeba zvolit přirozené číslo \(n\), abychom umocněním \( (1+x)^n \) dostali mnohočlen, ve kterém by měl koeficient kvadratického členu (koeficient u \( x^2 \)) hodnotu \( 300 \)?}
{\( 25 \)}{\( 24 \)}{\( 30 \)}{\( 20 \)}{}{}}
\LANGsk{
\Question{
Ako treba zvoliť prirodzené číslo \(n\), aby sme umocnením \( (1+x)^n \) dostali mnohočlen, v ktorom by mal koeficient kvadratického člena (koeficient pri \( x^2 \)) hodnotu \( 300 \)?}
{\( 25 \)}{\( 24 \)}{\( 30 \)}{\( 20 \)}{}{}}
\LANGes{
\Question{
Supón que el coeficiente del término cuadrático (es decir, el coeficiente en  \( x^2 \)) en la expansión de \( (1+x)^n \) es igual a  \( 300 \). Determina el valor del exponente \( n \).}
{\( 25 \)}{\( 24 \)}{\( 30 \)}{\( 20 \)}{}{}}
\End

\Data\ID{9000139710}
\Updated{2021-04-25 07:04:28}
\Level{C}
\SubArea{Combinatorics}
\LANGen{
\Question{
The wallet contains nine coins: three
\(1\)-Euro coins,
three \(2\)-Euro coins
and three \(5\)-Euro
coins. How many different amounts can be paid if we have to pay the amount exactly
and use just three coins for this payment?}
{\(\frac{5!}
{3!\, 2!}=10\)}{\(\frac{5!}
{3!}=20\)}{\(3^{3}=27\)}{\(3!=6\)}{}{}}
\LANGpl{
\Question{
W portfelu jest dziewięć monet: trzy monety o nominale
\(1\)-Euro,
trzy \(2\)-Euro 
i trzy \(5\)-Euro. Ile jest możliwości dokonania zapłaty, jeżeli musimy wydać  wyliczoną kwotę oraz użyć tylko trzech monet.}
{\(\frac{5!}
{3!\, 2!}=10\)}{\(\frac{5!}
{3!}=20\)}{\(3^{3}=27\)}{\(3!=6\)}{}{}}
\LANGcs{
\Question{
Kolik částek můžeme přesně zaplatit třemi mincemi? K dispozici máme
tři desetikorunové, tři dvacetikorunové a tři padesátikorunové mince.}
{\(\frac{5!}
{3!\, 2!}=10\)}{\(\frac{5!}
{3!}=20\)}{\(3^{3}=27\)}{\(3!=6\)}{}{}}
\LANGsk{
\Question{
V peňaženke máme deväť mincí: tri
\(1\)-eurové mince,
tri \(2\)-eurové mince
a tri \(5\)-eurové
mince. Koľko rôznych súm môžeme presne zaplatiť, ak na platbu použijeme len tri mince?}
{\(\frac{5!}
{3!\, 2!}=10\)}{\(\frac{5!}
{3!}=20\)}{\(3^{3}=27\)}{\(3!=6\)}{}{}}
\LANGes{
\Question{
Una cartera contiene nueve monedas: tres
monedas de \(1\) euro,
tres monedas de \(2\) euros 
y tres monedas de \(5\) euros. ¿Cuántas cantidades diferentes se pueden pagar si tenemos que pagar una cantidad exacta
y usar solo tres monedas?}
{\(\frac{5!}
{3!\, 2!}=10\)}{\(\frac{5!}
{3!}=20\)}{\(3^{3}=27\)}{\(3!=6\)}{}{}}
\End

\Data\ID{9000153901}
\Updated{2020-08-19 01:08:23}
\Level{C}
\SubArea{Combinatorics}
\LANGen{
\Question{
Find the number of ways how to distribute
\(8\) identical
balls among \(5\) 
persons so that each of them gets at least one ball.}
{\(\left({7\above 0.0pt 
3}\right) = 35\)}{\(5^{3} = 125\)}{\(\left({12\above 0.0pt 
 5} \right) = 792\)}{\(\left({12\above 0.0pt 
 8} \right) = 495\)}{}{}}
\LANGpl{
\Question{
Na ile sposobów można rozdzielić 
\(8\) identycznych
piłek wśród \(5\) osób tak, aby
każda z tych osób otrzymała przynajmniej jedną piłkę.}
{\(\left({7\above 0.0pt 
3}\right) = 35\)}{\(5^{3} = 125\)}{\(\left({12\above 0.0pt 
 5} \right) = 792\)}{\(\left({12\above 0.0pt 
 8} \right) = 495\)}{}{}}
\LANGcs{
\Question{
Otec má \(5\) 
synů a \(8\) 
stejných nerozlišitelných míčků. Kolika způsoby může synům míčky
rozdat, má-li každý dostat aspoň jeden?}
{\(\left ({7\above 0.0pt 
3}\right) = 35\)}{\(5^{3} = 125\)}{\(\left ({12\above 0.0pt 
 5} \right) = 792\)}{\(\left ({12\above 0.0pt 
 8} \right) = 495\)}{}{}}
\LANGsk{
\Question{
Koľkými spôsobmi je možné rozdať 
\(8\) rovnakých loptičiek \(5\) 
osobám, ak každý má dostať aspoň jednu loptičku?}
{\(\left({7\above 0.0pt 
3}\right) = 35\)}{\(5^{3} = 125\)}{\(\left({12\above 0.0pt 
 5} \right) = 792\)}{\(\left({12\above 0.0pt 
 8} \right) = 495\)}{}{}}
\LANGes{
\Question{
Determina el número de posibilidades de distribuir
\(8\) pelotas idénticas
entre \(5\)
personas para que cada uno obtenga al menos una pelota.}
{\(\left({7\above 0.0pt 
3}\right) = 35\)}{\(5^{3} = 125\)}{\(\left({12\above 0.0pt 
 5} \right) = 792\)}{\(\left({12\above 0.0pt 
 8} \right) = 495\)}{}{}}
\End

\Data\ID{9000153902}
\Updated{2020-08-19 01:08:22}
\Level{C}
\SubArea{Combinatorics}
\LANGen{
\Question{
Find the number of ways how to distribute
\(8\) identical
balls among \(5\) 
persons.}
{\(\left({12\above 0.0pt 
 8} \right) = 495\)}{\(5^{8} = 390625\)}{\(\left({8\above 0.0pt 
5}\right) = 56\)}{\(\left({12\above 0.0pt 
 5} \right) = 792\)}{}{}}
\LANGpl{
\Question{
Na ile sposobów można rozdzielić 
\(8\) identycznych
piłek wśród \(5\) 
osób.}
{\(\left({12\above 0.0pt 
 8} \right) = 495\)}{\(5^{8} = 390625\)}{\(\left({8\above 0.0pt 
5}\right) = 56\)}{\(\left({12\above 0.0pt 
 5} \right) = 792\)}{}{}}
\LANGcs{
\Question{
Otec má \(5\) 
synů a \(8\) 
stejných nerozlišitelných míčků. Kolika způsoby může míčky synům
rozdat?}
{\(\left ({12\above 0.0pt 
 8} \right) = 495\)}{\(5^{8} = 390625\)}{\(\left ({8\above 0.0pt 
5}\right) = 56\)}{\(\left ({12\above 0.0pt 
 5} \right) = 792\)}{}{}}
\LANGsk{
\Question{
Koľkými spôsobmi je možné rozdať
\(8\) rovnakých
loptičiek \(5\) 
osobám?}
{\(\left({12\above 0.0pt 
 8} \right) = 495\)}{\(5^{8} = 390625\)}{\(\left({8\above 0.0pt 
5}\right) = 56\)}{\(\left({12\above 0.0pt 
 5} \right) = 792\)}{}{}}
\LANGes{
\Question{
Determina el número de posibilidades de distribuir \(8\) pelotas idénticas entre \(5\) 
personas.}
{\(\left({12\above 0.0pt 
 8} \right) = 495\)}{\(5^{8} = 390625\)}{\(\left({8\above 0.0pt 
5}\right) = 56\)}{\(\left({12\above 0.0pt 
 5} \right) = 792\)}{}{}}
\End

\Data\ID{9000153903}
\Updated{2021-04-25 07:04:39}
\Level{C}
\SubArea{Combinatorics}
\LANGen{
\Question{
Find the number of ways how to distribute
\(8\) mutually different
balls among \(5\) 
persons.}
{\(5^{8} = 390625\)}{\(\frac{8!}
{3!} = 6720\)}{\(8^{5} = 32768\)}{\(\left({12\above 0.0pt 
 8} \right) = 495\)}{}{}}
\LANGpl{
\Question{
Na ile sposobów można rozdzielić
\(8\)  różnych piłek wśród  \(5\) 
osób.}
{\(5^{8} = 390625\)}{\(\frac{8!}
{3!} = 6720\)}{\(8^{5} = 32768\)}{\(\left({12\above 0.0pt 
 8} \right) = 495\)}{}{}}
\LANGcs{
\Question{
Otec má \(5\) 
synů a \(8\) 
různých míčků. Kolika způsoby může míčky synům rozdat?}
{\(5^{8} = 390625\)}{\(\frac{8!} 
{3!} = 6720\)}{\(8^{5} = 32768\)}{\(\left ({12\above 0.0pt 
 8} \right) = 495\)}{}{}}
\LANGsk{
\Question{
Koľkými spôsobmi je možné rozdať
\(8\) navzájom rôznych
loptičiek \(5\) 
osobám?}
{\(5^{8} = 390625\)}{\(\frac{8!}
{3!} = 6720\)}{\(8^{5} = 32768\)}{\(\left({12\above 0.0pt 
 8} \right) = 495\)}{}{}}
\LANGes{
\Question{
Determina el número de posibilidades de distribuir
\(8\)  pelotas diferentes 
entre \(5\) 
personas.}
{\(5^{8} = 390625\)}{\(\frac{8!}
{3!} = 6720\)}{\(8^{5} = 32768\)}{\(\left({12\above 0.0pt 
 8} \right) = 495\)}{}{}}
\End

\Data\ID{9000153905}
\Updated{2020-08-19 01:08:43}
\Level{C}
\SubArea{Combinatorics}
\LANGen{
\Question{
Find the number of ways how to distribute
\(5\) identical
balls among \(8\) 
persons.}
{\(\left({12\above 0.0pt 
 5} \right) = 792\)}{\(8^{5} = 32768\)}{\(\left({12\above 0.0pt 
 8} \right) = 495\)}{\(\frac{8!}
{3!} = 6720\)}{}{}}
\LANGpl{
\Question{
Na ile sposobów można rozdzielić 
\(5\) identycznych 
piłek wśród \(8\) 
osób.}
{\(\left({12\above 0.0pt 
 5} \right) = 792\)}{\(8^{5} = 32768\)}{\(\left({12\above 0.0pt 
 8} \right) = 495\)}{\(\frac{8!}
{3!} = 6720\)}{}{}}
\LANGcs{
\Question{
Otec má \(8\) 
synů a \(5\) 
stejných nerozlišitelných míčků. Kolika způsoby může míčky synům
rozdat?}
{\(\left ({12\above 0.0pt 
 5} \right) = 792\)}{\(8^{5} = 32768\)}{\(\left ({12\above 0.0pt 
 8} \right) = 495\)}{\(\frac{8!} 
{3!} = 6720\)}{}{}}
\LANGsk{
\Question{
Koľkými spôsobmi je možné rozdať
\(5\) rovnakých loptičiek
 \(8\) 
osobám?}
{\(\left({12\above 0.0pt 
 5} \right) = 792\)}{\(8^{5} = 32768\)}{\(\left({12\above 0.0pt 
 8} \right) = 495\)}{\(\frac{8!}
{3!} = 6720\)}{}{}}
\LANGes{
\Question{
Determina el número de posibilidades de distribuir 
\(5\) pelotas idénticas entre \(8\) 
personas.}
{\(\left({12\above 0.0pt 
 5} \right) = 792\)}{\(8^{5} = 32768\)}{\(\left({12\above 0.0pt 
 8} \right) = 495\)}{\(\frac{8!}
{3!} = 6720\)}{}{}}
\End

\Data\ID{9000153906}
\Updated{2021-04-25 07:04:11}
\Level{C}
\SubArea{Combinatorics}
\LANGen{
\Question{
Find the number of ways how to distribute
\(5\) identical
balls among \(8\) 
persons so that no persons gets more than one ball.}
{\(\frac{8!}
{5!3!} = 56\)}{\(\frac{8!}
{3!} = 6\:720\)}{\(\left({12\above 0.0pt 
 8} \right) = 495\)}{\(\left({12\above 0.0pt 
 5} \right) = 792\)}{}{}}
\LANGpl{
\Question{
Na ile sposobów można rozdzielić 
\(5\) identycznych piłek wśród  \(8\) 
osób, tak, aby żadna osoba nie otrzymała więcej niż jednej piłki.}
{\(\frac{8!}
{5!3!} = 56\)}{\(\frac{8!}
{3!} = 6\:720\)}{\(\left({12\above 0.0pt 
 8} \right) = 495\)}{\(\left({12\above 0.0pt 
 5} \right) = 792\)}{}{}}
\LANGcs{
\Question{
Otec má \(8\) 
synů a \(5\) 
stejných nerozlišitelných míčků. Kolika způsoby může míčky synům
rozdat, má-li každý dostat nejvýše jeden?}
{\(\frac{8!} 
{5!3!} = 56\)}{\(\frac{8!} 
{3!} = 6\:720\)}{\(\left ({12\above 0.0pt 
 8} \right) = 495\)}{\(\left ({12\above 0.0pt 
 5} \right) = 792\)}{}{}}
\LANGsk{
\Question{
Koľkými spôsobmi je možné rozdať
\(5\) rovnakých
loptičiek \(8\) 
osobám tak, aby každá osoba dostala najviac jednu loptičku?}
{\(\frac{8!}
{5!3!} = 56\)}{\(\frac{8!}
{3!} = 6\:720\)}{\(\left({12\above 0.0pt 
 8} \right) = 495\)}{\(\left({12\above 0.0pt 
 5} \right) = 792\)}{}{}}
\LANGes{
\Question{
Determina el número de posibilidades de distribuir 
\(5\) pelotas idénticas entre \(8\) 
personas para que ninguna persona obtenga más de una pelota.}
{\(\frac{8!}
{5!3!} = 56\)}{\(\frac{8!}
{3!} = 6\:720\)}{\(\left({12\above 0.0pt 
 8} \right) = 495\)}{\(\left({12\above 0.0pt 
 5} \right) = 792\)}{}{}}
\End

\Data\ID{9000153907}
\Updated{2020-08-19 01:08:31}
\Level{C}
\SubArea{Combinatorics}
\LANGen{
\Question{
Find the number of ways how to distribute
\(5\) different
balls among \(8\) 
persons.}
{\(8^{5} = 32\:768\)}{\(\left({12\above 0.0pt 
 5} \right) = 792\)}{\(\frac{8!}
{3!} = 6\:720\)}{\(\frac{8!}
{5!3!} = 56\)}{}{}}
\LANGpl{
\Question{
Na ile sposobów można rozdzielić 
\(5\) różnych piłek wśród \(8\) 
osób.}
{\(8^{5} = 32\:768\)}{\(\left({12\above 0.0pt 
 5} \right) = 792\)}{\(\frac{8!}
{3!} = 6\:720\)}{\(\frac{8!}
{5!3!} = 56\)}{}{}}
\LANGcs{
\Question{
Otec má \(8\) 
synů a \(5\) 
různých míčků. Kolika způsoby může synům míčky rozdat?}
{\(8^{5} = 32\:768\)}{\(\left ({12\above 0.0pt 
 5} \right) = 792\)}{\( \frac{8!} 
{3!} = 6\:720\)}{\(\frac{8!} 
{5!3!} = 56\)}{}{}}
\LANGsk{
\Question{
Koľkými spôsobmi je možné rozdať
\(5\) rôznych loptičiek \(8\) 
osobám?}
{\(8^{5} = 32\:768\)}{\(\left({12\above 0.0pt 
 5} \right) = 792\)}{\(\frac{8!}
{3!} = 6\:720\)}{\(\frac{8!}
{5!3!} = 56\)}{}{}}
\LANGes{
\Question{
Determina el número de posibilidades de distribuir
\(5\) pelotas diferentes entre \(8\) 
personas.}
{\(8^{5} = 32\:768\)}{\(\left({12\above 0.0pt 
 5} \right) = 792\)}{\(\frac{8!}
{3!} = 6\:720\)}{\(\frac{8!}
{5!3!} = 56\)}{}{}}
\End
